% !TeX root = ../../Book2.tex
% 纲要标题：## 第7章　投射模、内射模与平坦模
% 纲要位置：计划纲要.md，第 1272 行。

\chapter{投射模、内射模与平坦模}
\label{b2:ch:07}
\BookChapterNumber{b2:ch:07}

\begin{chapterabstract}
本章分别从同态的提升、延拓及张量后的正合性刻画投射模、内射模与平坦模，证明主要判据并比较三类性质。整数模、幂等元与多项式理想提供可直接检验的例子和反例，也说明自由、投射、平坦与无挠之间的蕴含不能任意倒转。
\end{chapterabstract}

\begin{learninggoals}
\begin{itemize}
\item 用提升与分裂刻画投射模，能够构造有限对偶基并识别不自由的投射模。
\item 理解内射性的延拓刻画与 Baer 判据，掌握极大延拓证明、可除交换群及足够多内射模的构造。
\item 区分自由、投射、平坦、无挠，使用有限关系或理想张量判据检验具体模。
\item 说明有限展示在平坦推出投射的证明中如何使用，并认识本章与后续局部理论、同调代数的关系。
\end{itemize}
\end{learninggoals}

对一个自由模，沿满射提升同态并不难：先为每个基向量选取原像，再线性延拓即可。若模有关系，这些原像还必须满足同样的关系，未必能选到。另一方面，把一个模与短正合列张量，满射仍会留下来，单射却可能消失；是否消失取决于所张量的模。这两类现象要求不同的条件，分别指向投射性和平坦性。还有一种方向相反的问题：定义在子模上的同态能否延伸到整个模，它将引出内射性。

\ProseParagraphBreak

三个问题对应于已经学过的两个 Hom 函子和张量函子的正合性。\(\operatorname{Hom}_R(P,-)\) 本来左正合，投射性要求它还保持右端满射；反变的 \(\operatorname{Hom}_R(-,E)\) 也本来左正合，内射性要求限制映射到子模后仍满射；\(-\otimes_RM\) 本来右正合，平坦性要求它还保持左端单射。因此三种性质分别补足一个函子通常可能失败的部分。以下判据将把这些对任意映射的要求化为可检验的模结构和关系条件。

\ProseParagraphBreak

本章需要\chapref{b2:ch:02}的模与同态、\chapref{b2:ch:03}的自由模、展示及分裂正合列，以及\chapref{b2:ch:06}的张量构造和右正合性。主理想整环上的例子另用\cref{b2:thm:pid-module-structure}。讨论平坦性时须将右模与左模配对。Baer 判据\footnote{贝尔（Reinhold Baer，1902-1979），德国裔美国数学家，研究群论及模的内射性质。}使用第一卷附录A定理A.1.4的 Zorn 引理。

\ProseParagraphBreak
\secref{b2:sec:0701}和\secref{b2:sec:0702}分别以提升和延拓刻画投射模、内射模；\secref{b2:sec:0703}用张量正合性定义平坦模，\secref{b2:sec:0704}则以反例检验自由、投射、平坦、无挠之间的逆向蕴含。第三卷的局部平坦性用到\cref{b2:thm:flat-ideal-criterion}，Morita 理论则用到\secref{b2:sec:0701}的有限生成投射模。

\ProseParagraphBreak
进一步阅读可参阅 Lam 的《Lectures on Modules and Rings》\cite[\S\S 2A--2B, 3A, 3C]{Lam1999ModulesRings}中的投射模、对偶基、Baer 判据及可除性；平坦性检验与有限展示条件见同书\cite[\S\S 4A, 4C--4E]{Lam1999ModulesRings}。不同文献采用的左模、右模约定可能不同。

% !TeX root = ../../Book2.tex
% 纲要标题：### 7.1 投射模
% 纲要位置：计划纲要.md，第 1274 行。

\section{投射模}
\label{b2:sec:0701}

若一个模有基，从它出发的同态便能沿满射提升：先为基向量的像选取原像，再线性延伸。没有基的模也可能具有这项性质。本节把它单独作为定义，并说明它恰好描述自由模的直和因子。除特别说明外，环均含幺元，模均为幺性左模。

% 纲要要点（供目录核对）：
%
% * 提升性质与自由模直和因子
% * 投射模和分裂正合列
% * 有限生成投射模的例子与幂等自同态的像；以有限对偶基就地证明 Hom 的张量实现
%

\subsection{提升性质与自由模直和因子}
\label{b2:subsec:0701:01}

\begin{definition}\label{b2:def:projective-module}
设 \(R\) 是环，\(P\) 是左 \(R\) 模。若对任意左 \(R\) 模之间的满同态 \(q:B\to C\) 和同态 \(f:P\to C\)，都存在同态 \(g:P\to B\) 满足 \(qg=f\)，就称 \(P\) 为\term[toushemo]{投射模}[projective module]。这里要求提升存在，不要求它唯一。
\end{definition}

这里要求下图中的虚箭头存在，并使两个从 \(P\) 到 \(C\) 的复合相等：
\[
\begin{tikzcd}[column sep=3em,row sep=2.5em]
& B \arrow[d,two heads,"q"] \\
P \arrow[r,"f"'] \arrow[ur,dashed,"g"] & C
\end{tikzcd}
\]

已知 \(f\) 的像在下面的 \(C\) 中；提升 \(g\) 则把这份映射沿满射抬到上面的 \(B\)，再由 \(q\) 映回时恢复 \(f\)。\cref{b2:prop:free-lifting}已经证明自由模具有这项性质。零模也投射。作为一个相反的例子，\(\mathbb Z/n\mathbb Z\) 在 \(n>1\) 时不是投射 \(\mathbb Z\) 模：若恒等映射能沿 \(\mathbb Z\to\mathbb Z/n\mathbb Z\) 提升，则 \(\bar1\) 的像应是被 \(n\) 零化的整数，只能为零，却又必须映回 \(\bar1\)。

\begin{theorem}\label{b2:thm:projective-characterizations}
设 \(R\) 是环，\(P\) 是左 \(R\) 模。下列条件等价：
\begin{equivalences}
\item\label{b2:item:projective-lifting} \(P\) 是投射模。
\item\label{b2:item:projective-split-epimorphism} 每个以 \(P\) 为终点的满同态都有模同态截面。
\item\label{b2:item:projective-summand} 存在左模 \(Q\)，使 \(P\oplus Q\) 是自由模。
\end{equivalences}
\end{theorem}
\begin{proof}
若 \(P\) 投射，对满同态 \(q:B\to P\) 提升 \(\operatorname{id}_P\)，便得到 \(qs=\operatorname{id}_P\)，即一个截面。

任意模都有自由满射 \(\pi:F\to P\)，例如取以 \(P\) 的底集合为基的自由模，把对应于 \(p\) 的基向量送到 \(p\)。若这个满射有截面，则\cref{b2:thm:split-exact-sequence}给出 \(F\cong\ker\pi\oplus P\)，所以第二项蕴含第三项。

最后设 \(F=P\oplus Q\) 自由，记包含与投影为 \(i:P\to F\)、\(p:F\to P\)，因而 \(pi=\operatorname{id}_P\)。给定定义中的 \(q\) 和 \(f\)，先用\cref{b2:prop:free-lifting}提升 \(fp:F\to C\)，得到 \(h:F\to B\) 且 \(qh=fp\)。令 \(g=hi\)，则 \(qg=fpi=f\)。
\end{proof}

定义要处理从 \(P\) 出发的所有提升问题，第二个条件把它缩成了以 \(P\) 为商模的分裂问题：每个短正合列 \(0\to A\to B\to P\to0\) 都应分裂。还可以只固定一个自由满射 \(\pi:F\to P\) 来检验；它有截面，就使 \(P\) 成为自由模的直和因子，因而投射。反过来，\(P\) 投射时，任何这样的自由满射都有截面。

\ProseParagraphBreak
这里“直和因子”包含嵌入与投影相容的意思。只有一个抽象单射 \(P\to F\)，并不足以证明 \(P\) 投射；\secref{b2:sec:0704}将给出嵌入自由模却不平坦、更不投射的理想。

\begin{proposition}\label{b2:prop:projective-sums-summands}
设 \(R\) 是环。左 \(R\) 投射模的任意直和是投射模，投射模的直和因子仍是投射模。
\end{proposition}
\begin{proof}
对 \(f:\bigoplus_iP_i\to C\) 和满射 \(q:B\to C\)，分别提升各限制 \(f_i:P_i\to C\)，得到 \(g_i:P_i\to B\)。\cref{b2:thm:module-sum-universal}将它们唯一合成 \(g:\bigoplus_iP_i\to B\)；\(qg=f\) 可在每个直和分量上验证。无限族的提升选择使用选择公理。

若 \(P\) 是投射模 \(T\) 的直和因子，用 \(i:P\to T\)、\(p:T\to P\) 且 \(pi=\operatorname{id}_P\) 表示它。先提升 \(fp\)，再与 \(i\) 复合，即得到 \(f\) 的提升。
\end{proof}

\subsection{投射模和分裂正合列}
\label{b2:subsec:0701:02}

投射性也可以用正合性表达。下面第一列的左端单射来自 \(A\to B\) 的单射，中间位置的正合性则说明：一个 \(P\to B\) 的同态落入 \(A\) 在 \(B\) 中的像，恰好等价于它复合到 \(C\) 为零。尚未保证的是每个 \(P\to C\) 都能提升到 \(B\)，这正是投射性要补足的末端满射。第二列将来源和目标的角色对换，下一节会用它把内射性解释为同态的延拓。

\begin{proposition}[Hom 的左正合性]\label{b2:prop:hom-left-exact}
设 \(R\) 是环，\(0\to A\xrightarrow{i}B\xrightarrow{q}C\to0\) 是左 \(R\) 模的短正合列，\(P,E\) 是左 \(R\) 模。下面两列交换群正合：
\[
0\longrightarrow\operatorname{Hom}_R(P,A)
\xrightarrow{i\circ-}\operatorname{Hom}_R(P,B)
\xrightarrow{q\circ-}\operatorname{Hom}_R(P,C),
\]
\[
0\longrightarrow\operatorname{Hom}_R(C,E)
\xrightarrow{-\circ q}\operatorname{Hom}_R(B,E)
\xrightarrow{-\circ i}\operatorname{Hom}_R(A,E).
\]
\end{proposition}
\begin{proof}
第一列中，\(i\) 单射保证 \(i\circ-\) 单射。若 \(qf=0\)，则 \(f(P)\subseteq\ker q=i(A)\)，由\cref{b2:prop:short-exact-kernel-cokernel}，\(f\) 唯一写成 \(ih\)，其中 \(h:P\to A\) 是同态。这给出中间位置的正合性。

第二列中，\(q\) 满射保证 \(-\circ q\) 单射。若 \(hi=0\)，则 \(h\) 在 \(i(A)\) 上为零，故公式 \(\bar h\bigl(q(b)\bigr)=h(b)\) 与代表元无关，并定义 \(\bar h:C\to E\)。它满足 \(h=\bar hq\)，且由 \(q\) 满射唯一确定。两列相邻映射的复合均为零，因此这些核与像的等式就是所需正合性。
\end{proof}

\begin{corollary}\label{b2:cor:projective-hom-exact}
设 \(R\) 是环，\(P\) 是左 \(R\) 模。\(P\) 投射，当且仅当函子 \(\operatorname{Hom}_R(P,-)\) 把每个左 \(R\) 模的短正合列变成交换群的短正合列；也等价于每个左 \(R\) 模短正合列 \(0\to A\to B\to P\to0\) 都分裂。
\end{corollary}
\begin{proof}
由\cref{b2:prop:hom-left-exact}，只需补上最后一个 Hom 映射的满射性，而它正是每个 \(P\to C\) 能沿 \(B\to C\) 提升的条件。任一满射与其核组成短正合列，所以这个条件覆盖定义中的所有满射。第二个等价性是\cref{b2:thm:projective-characterizations}与\cref{b2:thm:split-exact-sequence}的合用：满射有截面恰好使相应短正合列分裂。
\end{proof}

这个结论只要求右端的 \(P\) 投射。若 \(P\) 出现在左端，短正合列未必分裂。例如 \(0\to\mathbb Z\xrightarrow{\times n}\mathbb Z\to\mathbb Z/n\mathbb Z\to0\) 的左端自由，\(n>1\) 时仍不分裂。

\subsection{有限生成投射模的例子}
\label{b2:subsec:0701:03}

\begin{proposition}\label{b2:prop:finite-projective-summand}
设 \(R\) 是环，\(P\) 是左 \(R\) 模。\(P\) 有限生成且投射，当且仅当它是某个有限秩自由左 \(R\) 模 \(R^n\) 的直和因子；也等价于存在 \(n\ge0\) 和幂等自同态 \(E\in\operatorname{End}_R(R^n)\)，使 \(P\cong\operatorname{im}E\)。若 \(R\) 交换，后一条件可以用标准基写成幂等矩阵 \(E\in M_n(R)\)、\(E^2=E\) 的像。
\end{proposition}
\begin{proof}
若 \(P\) 有限生成，取满射 \(q:R^n\to P\)；若又投射，该满射有截面 \(s:P\to R^n\)，因而由\cref{b2:thm:split-exact-sequence}得到直和分解。反过来，有限自由模的直和因子由\cref{b2:prop:projective-sums-summands}知是投射模，而且是 \(R^n\) 在投影下的像，由标准基的有限个像生成。

任一这样的直和分解都有包含 \(s:P\to R^n\) 与投影 \(q:R^n\to P\)，满足 \(qs=\operatorname{id}_P\)。于是 \(E=sq\) 满足 \(E^2=s(qs)q=E\)，且 \(\operatorname{im}E=s(P)\cong P\)。反过来，若 \(E^2=E\)，每个 \(x\in R^n\) 都有分解
\[
x=E(x)+\bigl(x-E(x)\bigr),\qquad
E(x)\in\operatorname{im}E,\quad x-E(x)\in\ker E.
\]
若 \(y=E(x)\) 也在 \(\ker E\) 中，则 \(y=E^2(x)=E(y)=0\)，所以 \(R^n=\operatorname{im}E\oplus\ker E\)。因此 \(\operatorname{im}E\) 是有限自由模的直和因子。交换环上，标准基把自同态表示为通常作用于坐标列的矩阵，复合表示为矩阵乘法，故 \(E^2=E\) 正是幂等矩阵条件。
\end{proof}

设 \(e\in R\) 满足 \(e^2=e\)。作为左模，
\[
R=Re\oplus R(1-e).
\]
每个 \(r\) 都是 \(re+r(1-e)\)；若 \(x\) 同时属于两项，则右乘 \(e\) 分别得到 \(xe=x\) 与 \(xe=0\)，所以 \(x=0\)。因此 \(Re\) 是有限生成投射模。这一例子不要求 \(e\) 位于中心；右乘 \(e\) 才是这里的左模投影。

\ProseParagraphBreak
取 \(R=k\times k\)、\(e=(1,0)\)，则 \(Re=k\times0\) 非零，而 \((0,1)\) 零化的是整个模 \(Re\)，所以 \(\operatorname{Ann}_R(Re)\ne0\)。任意非零自由模却是忠实的：若 \(r\) 零化整个自由模，取其中一个基向量 \(e_i\)，则 \(re_i=0\) 的唯一坐标迫使 \(r=0\)。因此这个投射模不自由。这里必须比较整个模的零化子；仅有一个非零元素被非零标量零化，并不能排除自由性，正则自由模 \(R\) 中的 \(e\) 就被 \((0,1)\) 零化。

\ProseParagraphBreak
另一例是 \(R=M_n(k)\)、\(e=E_{11}\)：\(Re\) 由仅第一列可能非零的矩阵组成，可识别为通常的列向量模 \(k^n\)。当 \(n>1\) 时它也不自由，因为其 \(k\) 维数为 \(n\)，而非零自由 \(R\) 模若有有限基，其 \(k\) 维数为 \(n^2\) 的正整数倍。若有无限基，就可将它识别为 \(R^{(\Lambda)}\)，其中 \(\Lambda\) 无限。任意有限多个不同的标准基向量在 \(k\) 上也线性无关：在第 \(\lambda\) 个坐标读取一个线性关系，得到 \(a_\lambda I_n=0\)，故 \(a_\lambda=0\)。于是底层 \(k\) 空间的维数无限，与 \(\dim_k Re=n\) 矛盾。

\begin{proposition}[有限对偶基判据]\label{b2:prop:finite-dual-basis}
设 \(R\) 是环，\(P\) 是左 \(R\) 模。\(P\) 有限生成且投射，当且仅当存在有限个 \(p_1,\ldots,p_n\in P\) 和左 \(R\) 模同态 \(\varphi_i:P\to {}_RR\)，使对每个 \(p\in P\) 都有
\[
p=\sum_{i=1}^n\varphi_i(p)p_i.
\]
这样的元素与同态称为一组\term[youxian-duiouji]{有限对偶基}[finite dual basis]。它不要求 \((p_i)\) 线性无关，不要求 \(\varphi_i(p_j)=\delta_{ij}1_R\)，也不要求元素的系数表达唯一。
\end{proposition}
\begin{proof}
由\cref{b2:prop:finite-projective-summand}，有限生成投射模有分裂满射 \(q:R^n\to P\) 和截面 \(s:P\to R^n\)。令 \(p_i=q(e_i)\)，记第 \(i\) 个左坐标映射为 \(\pi_i:R^n\to R\)，并令 \(\varphi_i=\pi_i s\)。也就是说，先把 \(p\) 嵌回有限自由模，再读出第 \(i\) 个坐标。于是
\[
p=qs(p)=q\left(\sum_i\varphi_i(p)e_i\right)
=\sum_i\varphi_i(p)p_i.
\]

反过来，用 \(q(e_i)=p_i\) 定义 \(q:R^n\to P\)，用 \(s(p)=\sum_i\varphi_i(p)e_i\) 定义 \(s:P\to R^n\)。由于 \(\varphi_i(rp)=r\cdot\varphi_i(p)\)，\(s\) 确为左模同态。给定等式说明 \(qs=\operatorname{id}_P\)，所以 \(P\) 是有限自由模的直和因子。
\end{proof}

令 \(\psi_i:R\to P\)、\(r\mapsto rp_i\)，则重构公式也写成 \(\operatorname{id}_P=\sum_i\psi_i\varphi_i\)。每个复合都经过正则模 \(R\)，其像包含在循环子模 \(Rp_i\) 中；\(\varphi_i\) 和 \(\psi_i\) 都是左 \(R\) 线性的，因此这个解释在非交换环上也成立。有限对偶基是从有限自由模的直和分解继承的有限重构办法，允许 \((p_i)\) 有冗余。

\ProseParagraphBreak
沿用 \(R=k\times k\)、\(e=(1,0)\) 与 \(P=Re\) 的例子，取 \(p_1=e\)，并取包含映射 \(\varphi_1:P\hookrightarrow R\)。对 \(p=(a,0)\in P\)，有
\[
p=\varphi_1(p)e=(a,0)e,\qquad
\varphi_1(e)=e\ne1_R=(1,1).
\]
因此这一有限对偶基并不满足普通对偶基的 \(\varphi_i(p_j)=\delta_{ij}1_R\)。表示 \(p\) 的系数也不唯一：对每个 \(b\in k\)，\((a,b)e=(a,0)\)。有限对偶基给出一套指定的线性重构办法，却未必提供无关系的唯一坐标。

\ProseParagraphBreak
有限对偶基还把从 \(P\) 出发的 Hom 写成张量积。\cref{b2:prop:finite-free-hom-tensor}在有限自由情形用基与坐标泛函，把同态记录为有限多个基向量的像。对有限生成投射模，虽然 \((p_i)\) 不一定是基，有限对偶基仍使每个同态由有限多个值 \(f(p_i)\) 重构。正向对应仍把泛函 \(\lambda\) 与元素 \(m\) 组装为 \(p\mapsto\lambda(p)m\)；为了保留非交换环上的侧别，把额外的右作用一起写入陈述。

\begin{proposition}\label{b2:prop:finite-projective-hom-tensor}
设 \(R,S\) 为环，\({}_RP_S\) 为 \((R,S)\) 双模，且 \(P\) 作为左 \(R\) 模有限生成并投射。令 \(P^*=\operatorname{Hom}_R(P,{}_RR)\)，其 \((S,R)\) 双模作用为
\[
(s\lambda)(p)=\lambda(ps),\qquad
(\lambda r)(p)=\lambda(p)r.
\]
对每个左 \(R\) 模 \(M\)，令 \(F(M)=\operatorname{Hom}_R(P,M)\)，左 \(S\) 作用由 \((sf)(p)=f(ps)\) 给出；对左 \(R\) 模同态 \(a:M\to M'\)，令 \(F(a)(f)=af\)。则有关于 \(M\) 自然的左 \(S\) 模同构
\[
P^*\otimes_RM\longrightarrow F(M),\qquad
\lambda\otimes m\longmapsto\bigl(p\mapsto\lambda(p)m\bigr).
\]
\end{proposition}
\symidx{\(P^*=\operatorname{Hom}_R(P,R)\)：左模的右对偶模}
\begin{proof}
按式\eqref{b2:eq:hom-left-action}，\(P\) 的右 \(S\) 作用通过预复合，在 \(P^*\) 与 \(F(M)\) 上给出左 \(S\) 作用。\(P^*\) 的右 \(R\) 作用是在取值后右乘，也保留左 \(R\) 线性。两侧作用满足 \(\bigl((s\lambda)r\bigr)(p)=\lambda(ps)r=\bigl(s(\lambda r)\bigr)(p)\)，故 \(P^*\) 确为双模。由 \(\lambda(rp)=r\lambda(p)\)，\(p\mapsto\lambda(p)m\) 为左 \(R\) 线性。所写两变量映射双可加，且 \((\lambda r)(p)m=\lambda(p)(rm)\)，所以对 \(R\) 平衡；它也保留左 \(S\) 作用，因两种作用路径都在 \(p\) 处给出 \(\lambda(ps)m\)。故诱导所需同态。

由\cref{b2:prop:finite-dual-basis}，取有限个 \(p_i\in P,\lambda_i\in P^*\)，使 \(p=\sum_i\lambda_i(p)p_i\)。对任意同态 \(f:P\to M\)，左 \(R\) 线性给出 \(f(p)=\sum_i\lambda_i(p)f(p_i)\)。因此应把这些有限个值 \(f(p_i)\) 分别与重构泛函 \(\lambda_i\) 配对，反向定义 \(f\mapsto\sum_i\lambda_i\otimes f(p_i)\)。正向复合在 \(p\) 处取值为 \(\sum_i\lambda_i(p)f(p_i)=f(p)\)。另一复合将 \(\lambda\otimes m\) 送到
\[
\sum_i\lambda_i\otimes\lambda(p_i)m
=\left(\sum_i\lambda_i\cdot\lambda(p_i)\right)\otimes m
=\lambda\otimes m.
\]
最后一个等式来自对偶基公式施加 \(\lambda\)：\(\lambda(p)=\sum_i\lambda_i(p)\lambda(p_i)\)，其系数次序如式所写。故两映射互逆；正向映射左 \(S\) 线性，其逆也左 \(S\) 线性。对 \(M\) 上的同态，两种路径都将 \(\lambda\otimes m\) 送到 \(p\mapsto\lambda(p)a(m)\)，证明自然性。
\end{proof}

这个同构说明，有限生成投射模虽然未必有唯一的基坐标，有限对偶基仍足以把同态重构为有限张量和。在非交换环上，\(\lambda(p)\) 作用在 \(m\) 左侧，\(\lambda_i\cdot\lambda(p_i)\) 则使用 \(P^*\) 的右 \(R\) 作用；这些次序不能互换。“对偶基”这一名称并不允许忽略标量的作用方向。

% !TeX root = ../../Book2.tex
% 纲要标题：### 7.2 内射模
% 纲要位置：计划纲要.md，第 1280 行。

\section{内射模}
\label{b2:sec:0702}

投射模解决“沿满射提升”的问题。若给定的映射定义在一个子模上，希望把它延拓到整个模，所需的是另一种性质。向量空间可以选补空间后延拓线性映射，一般环上的模却未必有补模；内射性把这种延拓能力直接作为研究对象。

% 纲要要点（供目录核对）：
%
% * 延拓性质与 Baer 判据
% * 可除交换群、商群的可除性及内射中间项不保证分裂的例子
% * 模范畴中足够多内射对象
%

\subsection{延拓性质与 Baer 判据}
\label{b2:subsec:0702:01}

\begin{definition}\label{b2:def:injective-module}
设 \(R\) 是环，\(E\) 是左 \(R\) 模。若对任意左 \(R\) 模的子模包含 \(A\subseteq B\) 和同态 \(f:A\to E\)，都存在同态 \(g:B\to E\) 使 \(g|_A=f\)，就称 \(E\) 为\term[neishemo]{内射模}[injective module]。等价地，对任意左 \(R\) 模单同态 \(i:A\to B\) 及 \(f:A\to E\)，都能解出 \(gi=f\)。
\end{definition}

投射性固定同态的来源，沿满射作提升；内射性固定同态的目标 \(E\)，沿单射作延拓。定义中的延拓可画成下面的交换图；虚线是所求的同态，交换性即 \(gi=f\)。
\[
\begin{tikzcd}[column sep=3em,row sep=2.5em]
A\arrow[r,hook,"i"]\arrow[d,"f"']&B\arrow[dl,dashed,"g"]\\
E&
\end{tikzcd}
\]

\(\mathbb Z\) 不是内射 \(\mathbb Z\) 模：把 \(2\mathbb Z\) 映到 \(\mathbb Z\) 的同态 \(2n\mapsto n\) 若能延拓，则 \(2g(1)=1\)，在整数中无解。定义中的“内射”描述目标模的延拓性质，并不表示该模中的所有同态都是单射。

\begin{proposition}\label{b2:prop:injective-characterizations}
设 \(R\) 是环，\(E\) 是左 \(R\) 模。\(E\) 内射，等价于 \(\operatorname{Hom}_R(-,E)\) 将每个左 \(R\) 模的短正合列反向变成交换群的短正合列；也等价于每个左 \(R\) 模短正合列 \(0\to E\to B\to C\to0\) 都分裂。
\end{proposition}
\begin{proof}
\cref{b2:prop:hom-left-exact}说明 Hom 列只可能在最后一项不满，满射性恰是子模上的同态均能延拓。内射时，将恒等映射 \(E\to E\) 沿 \(E\to B\) 延拓，得到缩回，原列由\cref{b2:thm:split-exact-sequence}分裂。

反过来，设所有左端为 \(E\) 的短正合列分裂。给定左 \(R\) 模的子模包含 \(A\subseteq B\) 及同态 \(f:A\to E\)，希望把 \(E\) 和 \(B\) 放进同一个模，并强制 \(B\) 中的 \(a\) 与 \(E\) 中的 \(f(a)\) 相等。若 \(E\) 在这个模中有缩回，将 \(B\) 中的元素送回 \(E\) 就可得到延拓。为此，在 \(E\oplus B\) 中加入关系 \((f(a),-a)=0\)，按\cref{b2:prop:module-pullback-pushout}的推出构造取商模
\[
D=(E\oplus B)/\Bigl\{\,\bigl(f(a),-a\bigr):a\in A\,\Bigr\}.
\]
分母集合是同态 \(a\mapsto\bigl(f(a),-a\bigr)\) 的像，因而是子模。记商映射为 \(\pi:E\oplus B\to D\)。映射 \(j:E\to D\)、\(e\mapsto\pi(e,0)\) 单射：若 \((e,0)=\bigl(f(a),-a\bigr)\)，则 \(a=0,e=0\)。故 \(0\to E\xrightarrow{j}D\to D/j(E)\to0\) 有缩回 \(r:D\to E\)。令 \(g(b)=r\bigl(\pi(0,b)\bigr)\)，则对 \(a\in A\)，商中的等式 \(\pi(0,a)=\pi\bigl(f(a),0\bigr)\) 给出 \(g(a)=f(a)\)。
\end{proof}

这一构造把任意延拓问题改写为一个以 \(E\) 开头的分裂问题。与\cref{b2:thm:projective-characterizations}相比，投射性使所有以该模为右端的短正合列分裂，内射性则使所有以该模为左端的短正合列分裂。判断内射性时，还可以把需要检验的子模缩小到环自身的左理想：一次加入一个新元素时，它与旧定义域之间的全部关系恰由一个左理想记录。

\begin{theorem}[Baer 判据]\label{b2:thm:baer-criterion}
设 \(R\) 是含幺环，\(E\) 是左 \(R\) 模。\(E\) 内射，当且仅当对每个左理想 \(I\subseteq R\)，每个左 \(R\) 模同态 \(f:I\to E\) 都能延拓为左 \(R\) 模同态 \({}_RR\to E\)。后一条件也可写成：对每个这样的 \(I,f\)，存在 \(e\in E\)，使所有 \(a\in I\) 都满足 \(f(a)=ae\)。
\end{theorem}

这称为\term[baer-panju]{Baer 判据}[Baer's criterion]。等式形式来自 \(R\to E\) 的同态由 \(1\) 的像唯一确定。注意这里检验所有左理想；除非环有额外有限性条件，不能只检验有限生成的左理想。

\begin{proof}
必要性直接取 \(A=I,B=R\)。为证充分性，给定左 \(R\) 模的子模包含 \(A\subseteq B\) 和同态 \(f:A\to E\)，考虑所有延拓对 \((N,h)\)，其中 \(A\subseteq N\subseteq B\)、\(h:N\to E\)、\(h|_A=f\)。以定义域包含且映射限制相容为序。原始对 \((A,f)\) 保证集合非空；一条链的定义域之并仍是子模，相容的映射在其上拼成同态，给出上界；空链也有 \((A,f)\) 作上界。用第一卷附录A定理A.1.4的 Zorn 引理，取极大对 \((N,h)\)。

若 \(N\ne B\)，取 \(b\in B\setminus N\)，准备把 \(h\) 延拓到 \(N+Rb\)。若指定 \(h'(b)=e\)，则凡是 \(ab\in N\)，新旧映射都必须满足 \(h(ab)=ae\)。这些使新元素落回旧定义域的标量及旧映射强迫它们满足的条件，给出
\[
I=\{\,a\in R:ab\in N\,\},\qquad u:I\longrightarrow E,\quad u(a)=h(ab).
\]
\(I\) 对加法封闭，且 \(a\in I,r\in R\) 时，\((ra)b=r(ab)\in N\)，所以 \(I\) 是左理想；\(h\) 的线性又使 \(u\) 成为左模同态。按假设取 \(e\in E\) 使 \(u(a)=ae\)。一个延拓若把 \(b\) 送到 \(e\)，线性就强迫它只能按下式取值：
\[
h'(n+rb)=h(n)+re.
\]
所剩问题是这个公式会不会依赖 \(n+rb\) 的表示。若 \(n+rb=n'+r'b\)，则 \(a=r-r'\in I\) 且 \(ab=n'-n\)，所以 \(h(n'-n)=ae=(r-r')e\)，即
\[
h(n)+re=h(n')+r'e.
\]
因此 \(h'\) 良定义。加法与左标量相容由公式直接验证，且 \(h'|_N=h\)。它把极大延拓的定义域严格扩大，矛盾。因此 \(N=B\)，得到所需延拓。
\end{proof}

\subsection{可除交换群}
\label{b2:subsec:0702:02}

\begin{definition}\label{b2:def:divisible-group}
设 \(D\) 是交换群。若对每个整数 \(n>0\) 和每个 \(d\in D\)，都存在 \(x\in D\) 使 \(nx=d\)，就称 \(D\) 为\term[kechu-qun]{可除群}[divisible group]。这要求乘 \(n\) 为满射，不要求它为单射。
\end{definition}

\(\mathbb Q\) 与 \(\mathbb Q/\mathbb Z\) 都可除：分别用有理数除以 \(n\)，或把有理数代表元除以 \(n\)。后者有大量有限阶元素，说明可除与无挠是不同条件。非零有限交换群不可能可除，因为取群阶 \(n\)，乘 \(n\) 为零映射。

\ProseParagraphBreak
在整数环上，Baer 判据的延拓问题恰好变为这些除法方程。每个非零理想都是 \(n\mathbb Z\)，其中 \(n>0\)；给定 \(f:n\mathbb Z\to D\)，一个延拓由 \(x=g(1)\) 决定，所需条件只有 \(nx=f(n)\)。

\begin{theorem}\label{b2:thm:divisible-injective}
设 \(D\) 是交换群，并按整数倍看成左 \(\mathbb Z\) 模，则 \(D\) 内射，当且仅当它可除。若 \(D\) 可除，则对每个子群 \(A\leq D\)，商群 \(D/A\) 也可除，因而内射。但是交换群的短正合列
\[
0\longrightarrow\mathbb Z\longrightarrow\mathbb Q
\longrightarrow\mathbb Q/\mathbb Z\longrightarrow0,
\]
其中箭头分别为包含及自然商映射，并不分裂；它的中间项和右端均为内射 \(\mathbb Z\) 模。
\end{theorem}
\begin{proof}
若 \(D\) 内射，给定 \(d\in D\) 与 \(n>0\)，同态 \(n\mathbb Z\to D\)、\(nk\mapsto kd\) 可延拓为 \(g:\mathbb Z\to D\)。取 \(x=g(1)\)，便有 \(nx=g(n)=d\)。

若 \(D\) 可除，用\cref{b2:thm:baer-criterion}检验整数理想。零理想的映射用零映射延拓；对 \(n\mathbb Z\)、\(n>0\)，给定 \(f\)，取 \(x\) 使 \(nx=f(n)\)。令 \(g(k)=kx\)，则 \(g(nk)=k\cdot f(n)=f(nk)\)，得到延拓。

再设 \(D\) 可除，\(A\leq D\)。给定 \(d+A\in D/A\) 和正整数 \(n\)，由可除性取 \(x\in D\) 使 \(nx=d\)，则 \(n(x+A)=d+A\)。所以商群可除，由刚证得的等价性又知它内射。特别地，\(\mathbb Q\) 可除，其商 \(\mathbb Q/\mathbb Z\) 也可除，两者都内射。

若所写短正合列分裂，\cref{b2:thm:split-exact-sequence}给出缩回 \(\rho:\mathbb Q\to\mathbb Z\)，且 \(\rho(1)=1\)。但群同态又要求 \(2\rho(1/2)=\rho(1)=1\)，在整数中不可能。因此该列不分裂。这不违背内射模的分裂刻画：它的左端是并不内射的 \(\mathbb Z\)。
\end{proof}

\(\mathbb Z\) 投射但不内射。在域上，每个向量空间都有基和补空间，所以每个模既投射又内射；离开域以后，这两个条件分别约束同态的出发端和到达端。

\subsection{模范畴中足够多内射对象}
\label{b2:subsec:0702:03}

每个左 \(R\) 模 \(M\) 都有自由模的满射 \(F\twoheadrightarrow M\)，自由模又投射，所以模范畴有足够多投射对象。内射方向需要证明相反的存在性：总能找到内射模 \(E\) 及单射 \(M\hookrightarrow E\)。

\ProseParagraphBreak
先从底交换群着手。若要用映入某个交换群的同态检测非零元素，\(\mathbb Q\) 无法检测有限阶元素；\(\mathbb Q/\mathbb Z\) 却含有每种有限阶所需的像，又是内射交换群，可以把循环子群上的同态延拓到整个群。这给出下述分离性质。

\begin{lemma}\label{b2:lem:qz-separates}
设 \(A\) 是交换群，\(0\ne a\in A\)。存在群同态 \(\chi:A\to\mathbb Q/\mathbb Z\) 使 \(\chi(a)\ne0\)。
\end{lemma}
\begin{proof}
先在 \(a\) 生成的循环子群上定义映射。若 \(a\) 的阶为 \(n>1\)，令 \(a\mapsto 1/n+\mathbb Z\)；若 \(a\) 为无限阶，令 \(a\mapsto1/2+\mathbb Z\)。两种情况均给出循环群同态，且 \(a\) 的像非零。\cref{b2:thm:divisible-injective}说明 \(\mathbb Q/\mathbb Z\) 内射，所以该同态可延拓到 \(A\)。
\end{proof}

然后把交换群上的延拓转回 \(R\) 模上的延拓。对保单位元的同态 \(\mathbb Z\to R\)，限制标量就是取底交换群；\cref{b2:prop:scalar-change-adjunctions}给出的右伴随是余诱导 \(D\mapsto\operatorname{Hom}_{\mathbb Z}(R,D)\)。它将映入 \(D\) 的群同态与映入余诱导模的 \(R\) 线性映射相对应，因此可用 \(D\) 的内射性构造一个内射 \(R\) 模。

\begin{lemma}\label{b2:lem:coinduced-injective}
设 \(R\) 是含幺环，\(D\) 是作为 \(\mathbb Z\) 模内射的交换群，则
\[
J=\operatorname{Hom}_{\mathbb Z}(R,D),\qquad (rf)(s)=f(sr)
\]
是内射左 \(R\) 模。
\end{lemma}
\begin{proof}
这里 \(R\) 在 Hom 中先作为加法群。由\cref{b2:prop:scalar-change-adjunctions}，所写作用使 \(J\) 成为左 \(R\) 模；对任意左 \(R\) 模 \(M\)，其伴随同构给出互逆对应
\[
\begin{aligned}
\operatorname{Hom}_R(M,J)&\longleftrightarrow\operatorname{Hom}_{\mathbb Z}(M,D),\\
\alpha&\longmapsto\bigl(m\mapsto\alpha(m)(1)\bigr),\\
\chi&\longmapsto\Bigl[m\mapsto\bigl(s\mapsto\chi(sm)\bigr)\Bigr].
\end{aligned}
\]
对给定的 \(\chi\)，记第二个公式所得的左 \(R\) 线性映射为 \(\alpha_\chi\)，即 \(\alpha_\chi(m)(s)=\chi(sm)\)。

若 \(A\subseteq B\) 是左 \(R\) 模的子模包含，给定 \(A\to J\)，把它变为 \(A\to D\) 的加法群同态 \(\chi\)。由 \(D\) 内射，延拓为 \(\widetilde\chi:B\to D\)，再用反向公式得到 \(\alpha_{\widetilde\chi}:B\to J\)。在 \(a\in A,s\in R\) 处有 \(\widetilde\chi(sa)=\chi(sa)\)，所以后者确实延拓原映射。
\end{proof}

\begin{theorem}[足够多内射对象]\label{b2:thm:enough-injectives}
对任意含幺环 \(R\)，每个左 \(R\) 模 \(M\) 都能嵌入一个内射左模。
\end{theorem}
\begin{proof}
令 \(D=\mathbb Q/\mathbb Z\)、\(J=\operatorname{Hom}_{\mathbb Z}(R,D)\)。先对每个底交换群同态 \(\chi:M\to D\)，按\cref{b2:lem:coinduced-injective}中的伴随对应取左 \(R\) 线性映射 \(\alpha_\chi:M\to J\)，其中 \(\alpha_\chi(m)(s)=\chi(sm)\)。一个这样的映射未必检测所有非零元素；于是以集合 \(H=\operatorname{Hom}_{\mathbb Z}(M,D)\) 为指标，将它们全部组成积映射
\[
\iota:M\longrightarrow\prod_{\chi\in H}J,
\qquad \iota(m)=\bigl(\alpha_\chi(m)\bigr)_{\chi\in H}.
\]
积的泛性质保证 \(\iota\) 左线性。若 \(m\ne0\)，\cref{b2:lem:qz-separates}给出 \(\chi(m)\ne0\) 的分量；在 \(s=1\) 处有 \(\alpha_\chi(m)(1)=\chi(m)\ne0\)，故 \(\iota\) 单射。

由\cref{b2:thm:divisible-injective}及\cref{b2:lem:coinduced-injective}，每个 \(J\) 内射，任意积 \(\prod J\) 也内射：对子模上的映射逐分量延拓，再用\cref{b2:thm:module-product-universal}合成一个映入积的同态。这种合成不要求有限支集；分量族无限时，选择各延拓使用选择公理。因此所构造的积就是所需内射模。
\end{proof}

内射模的积封闭性与\cref{b2:prop:projective-sums-summands}中的投射模直和封闭性相对应：映入积的同态按坐标延拓，从直和出发的同态按分量提升，分别使用积与余积的泛性质。

\ProseParagraphBreak
“\term[zugouduo-neisheduixiang]{足够多内射对象}[enough injectives]”意指上述嵌入总能找到；不意指每个模本身内射，也不意指内射模可以选得有限生成。例如，有限生成可除交换群只能为零：按\cref{b2:thm:pid-module-structure}分成有限秩自由部分与有限挠部分后，乘二的满射性排除非零自由部分，再乘挠部分的阶便排除非零挠部分。因此任何包含 \(\mathbb Z\) 的内射交换群都不可能有限生成。反复对商模作这样的嵌入，将在第四卷构造内射消解。

% !TeX root = ../../Book2.tex
% 纲要标题：### 7.3 平坦模
% 纲要位置：计划纲要.md，第 1286 行。

\section{平坦模}
\label{b2:sec:0703}

\cref{b2:thm:tensor-right-exact}说明张量积总保持余核，却未必保持单射。最简单的失败已经出现在整数模中：取 \(n>1\)，把 \(\mathbb Z\xrightarrow{\times n}\mathbb Z\) 与 \(\mathbb Z/n\mathbb Z\) 张量，所得映射为零。平坦性正是要求这种失败不发生。本节先讨论一般环；理想张量判据及其后续应用处再明确假设环交换。

% 纲要要点（供目录核对）：
%
% * 张量正合性与平坦性
% * 自由、投射、平坦之间的蕴含
% * 主理想整环上的无挠与平坦
% * 在交换环情形就地证明有限生成理想张量判据：\(M\) 平坦当且仅当每个有限生成理想 \(I\) 的映射 \(I\otimes_RM\to M\) 都是单射；采用有限关系判据组织证明
% * 双数环的平坦判据及平方零算子的 Jordan 块解释
% * 第三卷第12.1节在此基础上建立局部判据，再由第三卷第13.4节用于完备化；本节承担一般理想张量判据，交换代数卷承担局部检测与应用
%

\subsection{张量正合性与平坦性}
\label{b2:subsec:0703:01}

\begin{definition}\label{b2:def:flat-module}
设 \(R\) 是含幺环，\(M\) 是幺性左 \(R\) 模。若对每个幺性右 \(R\) 模的单同态 \(u:A\to B\)，诱导的交换群同态
\[
u\otimes\operatorname{id}_M:A\otimes_R M\longrightarrow B\otimes_R M
\]
仍为单射，就称 \(M\) 为\term[pingtanmo]{平坦模}[flat module]。右模平坦的定义相应使用它与左模的张量积，不能在非交换环上省略这一侧别。
\end{definition}

对包含 \(A\hookrightarrow B\)，平坦性要求 \(A\otimes_RM\) 中本来不同的两个元素，进入 \(B\otimes_RM\) 后仍不同。等价地，一个由 \(A\) 中元素写成的张量有限和，若借助 \(B\) 中更多元素和张量关系成为零，它在 \(A\otimes_RM\) 中就已经为零。后面的有限关系判据将把这种“不产生新的零关系”写成系数的具体条件。

\begin{proposition}\label{b2:prop:flat-exact}
设 \(R\) 是含幺环，\(M\) 是幺性左 \(R\) 模。\(M\) 平坦，当且仅当 \(-\otimes_R M\) 将每个幺性右 \(R\) 模短正合列变成交换群的短正合列。
\end{proposition}
\begin{proof}
对 \(0\to A\to B\to C\to0\)，\cref{b2:thm:tensor-right-exact}已经给出 \(A\otimes_R M\to B\otimes_R M\to C\otimes_R M\to0\) 的正合性。平坦性补上左端单射。反过来，任意单射与其余核组成短正合列，所以保持这些短正合列便保证保持所有单射。
\end{proof}

张量所得通常只是交换群的正合列。若 \(R\) 交换，可以用剩余的标量作用把它们看成 \(R\) 模；若存在额外双模结构，也须先说明该作用是什么。

\begin{proposition}\label{b2:prop:flat-torsionfree}
设 \(R\) 是整环，\(M\) 是幺性平坦 \(R\) 模，则 \(M\) 无挠。
\end{proposition}
\begin{proof}
设 \(0\ne a\in R\)。整环上乘 \(a\) 给出单射 \(R\to R\)，与 \(M\) 张量后，经单位同构 \(R\otimes_R M\cong M\) 化为乘 \(a\) 的映射 \(M\to M\)。由平坦性它单射，因此 \(am=0\) 只能有 \(m=0\)。
\end{proof}

平坦性的定义比较的是任意右模的单射；整环上无挠只检验上述一类乘法单射。把这一个必要条件误作充分条件，正是一般整环上容易出现的错误。

\subsection{自由模、投射模与平坦模}
\label{b2:subsec:0703:02}

\cref{b2:prop:tensor-direct-sums}已经证明，固定右模 \(N\) 时有自然同构
\[
N\otimes_R\left(\bigoplus_iM_i\right)
\cong\bigoplus_i(N\otimes_R M_i).
\]
它把 \(n\otimes(m_i)_i\) 送到 \((n\otimes m_i)_i\)，并与 \(N\) 的同态相容。因此单射在张量直和后的表现，可以逐分量判断。

\begin{theorem}\label{b2:thm:projective-flat}
设 \(R\) 是含幺环，所论左 \(R\) 模均为幺性模。自由左 \(R\) 模是平坦模；平坦左 \(R\) 模的任意直和及直和因子仍平坦。因此每个投射左 \(R\) 模都平坦。
\end{theorem}
\begin{proof}
对 \(M=R\)，单位同构把 \(u\otimes\operatorname{id}_R\) 识别为原单射 \(u\)。对自由模 \(R^{(I)}\)，\cref{b2:prop:tensor-direct-sums}把它识别为 \(u\) 的若干份直和，仍为单射。同理，若每个 \(M_i\) 平坦，则相应诱导映射逐分量单射，它们的直和也单射。

若 \(M\oplus M'\) 平坦，则 \(u\otimes\operatorname{id}_{M\oplus M'}\) 是两个分量映射的直和。一个直和映射单射时，每个分量均单射，所以 \(M\) 平坦。最后用\cref{b2:thm:projective-characterizations}：投射模是自由模的直和因子。
\end{proof}

于是对任意环，已有
\[
\text{自由}\Longrightarrow\text{投射}\Longrightarrow\text{平坦}.
\]
整环上还可接上“平坦 \(\Longrightarrow\) 无挠”。在域上这些条件完全相同；在主理想整环上，对有限生成模，它们也由\cref{b2:thm:pid-module-structure}全部等价。在一般环上，投射不必自由的例子已由 \((k\times k)(1,0)\) 给出；下面的 \(\mathbb Q\) 将说明平坦也不必投射。

\subsection{主理想整环上的无挠与平坦}
\label{b2:subsec:0703:03}

有限生成无挠的主理想整环模自由，已由\cref{b2:thm:pid-module-structure}得到。任意无挠模未必有限生成，但其中每次张量计算只涉及有限多个元素和有限条关系，这使有限生成的结论仍然能够使用。

\begin{lemma}\label{b2:lem:tensor-directed-union}
设 \(R\) 是含幺环，幺性左 \(R\) 模 \(M\) 是一族子模 \((M_\lambda)_{\lambda\in\Lambda}\) 的并，且任意两项都包含在某个共同的较大项中。若每个 \(M_\lambda\) 平坦，则 \(M\) 平坦。
\end{lemma}
\begin{proof}
张量元素的表达只用有限多个元素，证明这个表达为零也只用有限多个定义关系。因此表达和零等式的证明都能落在一个较大的子模中。具体地，对任意右 \(R\) 模 \(N\)，\(z\in N\otimes_RM\) 是有限个纯张量之和。为其中每个 \(M\) 元素各选一个包含它的子模，再反复使用两个子模有共同较大项的条件，得到一个包含全部这些元素的 \(M_\lambda\)，所以 \(z\) 来自 \(N\otimes_RM_\lambda\)。

若某个 \(z_\lambda\in N\otimes_RM_\lambda\) 的像在 \(N\otimes_RM\) 中为零，根据张量积的构造，它在自由交换群中的代表元属于定义关系生成的子群。属于这个子群意味着它是有限个双可加关系和平衡关系的整数线性组合；这就是零等式的一个有限证明。为这些关系中出现的有限多个 \(M\) 元素各选一个包含它的子模，再与 \(M_\lambda\) 一起取共同较大项 \(M_\mu\)。相同的关系已经证明 \(z_\lambda\) 的像在 \(N\otimes_RM_\mu\) 中为零。这里没有假定 \(N\otimes_RM_\lambda\to N\otimes_RM\) 事先单射。

现在给定右模单射 \(u:A\to B\)。若 \(z\in A\otimes_R M\) 映到零，先在某个 \(A\otimes_R M_\lambda\) 表示它，再由上一段选取较大的 \(M_\mu\)，使其像已在 \(B\otimes_R M_\mu\) 中为零。\(M_\mu\) 平坦，所以所选提升在 \(A\otimes_R M_\mu\) 中为零，进而 \(z\) 在 \(A\otimes_R M\) 中为零。这就证明了单射性。
\end{proof}

\begin{theorem}\label{b2:thm:pid-flat-torsionfree}
设 \(R\) 是主理想整环，\(M\) 是幺性 \(R\) 模。\(M\) 平坦，当且仅当它无挠；这里不要求有限生成。
\end{theorem}
\begin{proof}
必要性是\cref{b2:prop:flat-torsionfree}。若 \(M\) 无挠，则每个有限生成子模 \(N\) 都无挠，由\cref{b2:thm:pid-module-structure}自由，因而平坦。全部有限生成子模之并为 \(M\)，两个这样的子模包含在它们的和中，后者仍有限生成。因此可应用\cref{b2:lem:tensor-directed-union}。
\end{proof}

这个证明使用有限生成条件的位置，是每个子模 \(N\) 的自由性，而不是整个 \(M\) 的大小。任意具体的张量表达及零关系，都能在某个这样的自由子模中处理，再由有向并传回 \(M\)。

\ProseParagraphBreak
对含幺环 \(R\)，非零投射左 \(R\) 模 \(P\) 必有一个非零同态 \(P\to R\)：由\cref{b2:thm:projective-characterizations}，它是某个自由模 \(R^{(I)}\) 的直和因子，因而有单射 \(i:P\to R^{(I)}\)。取非零 \(p\in P\)，则 \(i(p)\ne0\)，必有一个坐标非零，所以相应坐标投影与 \(i\) 的复合非零。因此，非零模 \(P\) 若满足 \(\operatorname{Hom}_R(P,R)=0\)，就不可能投射。

\ProseParagraphBreak
例如 \(\mathbb Q\) 无挠，由上一定理是平坦 \(\mathbb Z\) 模，却不是投射模。事实上，任一 \(h:\mathbb Q\to\mathbb Z\) 都为零：对 \(q\in\mathbb Q\) 和每个正整数 \(n\)，有 \(h(q)=n\cdot h(q/n)\)，故 \(h(q)\) 被所有正整数整除，只能为零。它是非零模，却没有非零同态到 \(\mathbb Z\)，所以不投射。

\subsection{有限关系与理想张量判据}
\label{b2:subsec:0703:04}

要有效检验平坦性，不能逐一列举所有模之间的单射。下面的判据把它变为对有限线性关系的要求。先保留一般环，并始终把系数写在左模元素的左侧。

\begin{theorem}[有限关系判据]\label{b2:thm:flat-finite-relations}
设 \(R\) 是含幺环，\(M\) 是幺性左 \(R\) 模。\(M\) 平坦，当且仅当每个有限关系
\[
\sum_{i=1}^n a_i x_i=0\qquad(a_i\in R,\ x_i\in M)
\]
都可写成如下形式：存在有限个 \(y_1,\ldots,y_t\in M\) 与 \(b_{ij}\in R\)，使
\begin{equation}\label{b2:eq:flat-relation-factorization}
x_i=\sum_{j=1}^t b_{ij}y_j,\qquad
\sum_{i=1}^n a_i b_{ij}=0\quad(1\le j\le t).
\end{equation}
允许 \(t=0\)，此时所有 \(x_i\) 均为零。
\end{theorem}

\term[youxian-guanxi-panju]{有限关系判据}[equational criterion for flatness]说明：模中的一个关系，可以由系数环中已经为零的有限关系产生。记 \(A=(a_1,\ldots,a_n)\)、\(X=(x_i)^{\mathsf T}\)、\(B=(b_{ij})\)、\(Y=(y_j)^{\mathsf T}\)，则判据中的条件为
\[
AX=0\quad\Longrightarrow\quad X=BY\ \text{且}\ AB=0.
\]
这里 \(AX\)、\(BY\) 编码系数在左侧作用的有限求和，\(AB\) 编码按 \(a_i b_{ij}\) 次序计算的乘积。在非交换环上，这种记法并不把 \(X\mapsto AX\) 当作左 \(R\) 线性映射，转置也不能改变乘法次序。判据没有要求 \(x_i\) 本身线性无关，也不要求找到统一适用于所有关系的一组 \(y_j\)。

\begin{proof}
设 \(M\) 平坦。要找到满足 \(\sum_i a_i b_{ij}=0\) 的系数列，就把所有这样的列组织成一个核。取右自由模映射 \(\rho:R_R^n\to R_R\)，\(\rho(e_i)=a_i\)，令 \(K=\ker\rho\)、\(I=\operatorname{im}\rho\)。一个向量 \(\sum_i e_i b_i\) 属于 \(K\)，恰好意味着 \(\sum_i a_i b_i=0\)，所以 \(K\) 记录系数行的全部关系。短正合列 \(0\to K\to R^n\to I\to0\) 张量后正合，而 \(I\otimes_RM\to R\otimes_RM\) 单射。因此 \((x_i)\in M^n\) 在 \(\rho\otimes\operatorname{id}_M\) 下为零时，来自 \(K\otimes_RM\) 中的某个有限和 \(\sum_j k_j\otimes y_j\)。写 \(k_j=\sum_i e_i b_{ij}\)，便得到 \eqref{b2:eq:flat-relation-factorization}，其中第二个等式正是 \(k_j\in K\)。

反过来，设所述单个关系都能这样分解。自由模中的一个张量为零，意味着它的每个自由坐标都为零；要同时处理这些坐标，就须让有限多个关系共用同一个分解。若 \(a_{\ell i}\in R\) 与 \(x_i\in M\) 满足 \(\sum_i a_{\ell i}x_i=0\)（\(1\le\ell\le s\)），则存在有限个 \(y_j\in M\) 与 \(b_{ij}\in R\)，使 \(x_i=\sum_jb_{ij}y_j\) 且 \(\sum_i a_{\ell i}b_{ij}=0\) 对所有 \(\ell,j\) 成立。对关系的条数归纳。零条时取 \(y_i=x_i\)、\(b_{ij}=\delta_{ij}\)；一条时为假设。多条时先用第一条得到 \(x_i=\sum_jb^{(1)}_{ij}y_j\)，且 \(\sum_i a_{1i}b^{(1)}_{ij}=0\)。将其余关系代入，得到关于 \(y_j\) 的较少条关系，再用归纳假设写 \(y_j=\sum_kb^{(2)}_{jk}z_k\)。取 \(b_{ik}=\sum_jb^{(1)}_{ij}b^{(2)}_{jk}\)，所有关系便同时为零。此处保留乘法顺序，没有使用交换性。

有了同时分解，就能检验自由模中任意子模的包含。设 \(K\subseteq F\)，其中 \(F\) 为任意右自由模。若 \(z=\sum_{i=1}^n k_i\otimes x_i\) 映到 \(F\otimes_RM\) 后为零，每个 \(k_i\) 只有有限支集，有限多个这样的支集之并仍有限。取覆盖这个并的基向量 \(e_1,\ldots,e_s\)，写 \(k_i=\sum_\ell e_\ell a_{\ell i}\)。由\cref{b2:prop:tensor-free-coordinate}，在 \(F\otimes_RM\cong\bigoplus M\) 的坐标中，条件就是 \(\sum_i a_{\ell i}x_i=0\) 对所有 \(\ell\) 成立。由上一段取 \(x_i=\sum_jb_{ij}y_j\)，且 \(\sum_i a_{\ell i}b_{ij}=0\) 对所有 \(\ell,j\) 成立。于是
\[
z=\sum_j\left(\sum_i k_i b_{ij}\right)\otimes y_j=0,
\]
因为每个内和的所有自由坐标均为零。故 \(K\otimes_R M\to F\otimes_R M\) 单射。

任意右模包含 \(U\subseteq V\) 未必直接位于一个自由模中，但可以先用自由满射 \(p:F\to V\) 覆盖它。取原像 \(K=p^{-1}(U)\)，则 \(K\) 是自由模 \(F\) 的子模，而限制 \(\pi:K\to U\) 仍满射：每个 \(u\in U\) 的 \(p\) 原像自动落在 \(K\)。这使 \(U\otimes_RM\) 中的元素可以提升到上一段能处理的子模。令 \(L=\ker p\)，记 \(\ell:L\to K\)、\(\iota:K\to F\) 与 \(j:U\to V\) 为包含映射。它们组成下图，两行均为短正合列：
\[
\begin{tikzcd}[column sep=2.2em,row sep=2.5em]
0\arrow[r]&L\arrow[r,hook,"\ell"]\arrow[d,equal]&K\arrow[r,"\pi"]\arrow[d,hook,"\iota"]&U\arrow[r]\arrow[d,hook,"j"]&0\\
0\arrow[r]&L\arrow[r,hook,"\iota\ell"']&F\arrow[r,"p"']&V\arrow[r]&0.
\end{tikzcd}
\]
若 \(z\in U\otimes_R M\) 在 \(V\otimes_R M\) 中的像为零，上排的张量右正合性保证 \(\pi\otimes\operatorname{id}_M\) 满射，所以可选 \(w\in K\otimes_R M\)，使 \((\pi\otimes\operatorname{id}_M)(w)=z\)。记它在 \(F\otimes_R M\) 中的像为 \(w_F\)。由图的交换性，\((p\otimes\operatorname{id}_M)(w_F)=(j\otimes\operatorname{id}_M)(z)=0\)。下排的张量右正合性又说明，\(w_F\) 来自 \(L\otimes_R M\)，故可选 \(v\in L\otimes_R M\)，使
\[
w_F=((\iota\ell)\otimes\operatorname{id}_M)(v).
\]

现在令 \(w'=w-(\ell\otimes\operatorname{id}_M)(v)\)。如同\cref{b2:thm:snake-lemma}证明中对提升的修正，减去的项测量并消去所选提升的误差。这个修正同时完成两个目标：因为 \(\pi\ell=0\)，它不改变在 \(U\otimes_RM\) 中的像；因为减去的正是 \(w_F\) 的原像，它使在 \(F\otimes_RM\) 中的像变成零。具体地，
\[
\begin{aligned}
(\pi\otimes\operatorname{id}_M)(w')&=z,\\
(\iota\otimes\operatorname{id}_M)(w')&=w_F-((\iota\ell)\otimes\operatorname{id}_M)(v)=0.
\end{aligned}
\]
以上提升与修正只用了张量右正合性。最后使用上一段已经证明的 \(K\subseteq F\) 的张量单射性，得到 \(w'=0\)，于是 \(z=(\pi\otimes\operatorname{id}_M)(w')=0\)。这证明 \(M\) 平坦。
\end{proof}

\begin{corollary}[有限关系组分解]\label{b2:cor:flat-finite-relation-systems}
设 \(R\) 是含幺环，\(M\) 是幺性平坦左 \(R\) 模。若有限个 \(x_i\in M\) 对 \(a_{\ell i}\in R\) 满足
\[
\sum_{i=1}^n a_{\ell i}x_i=0\qquad(1\le\ell\le s),
\]
则存在有限个 \(y_j\in M\) 与 \(b_{ij}\in R\)，使
\[
x_i=\sum_j b_{ij}y_j,\qquad
\sum_i a_{\ell i}b_{ij}=0\quad\text{对所有 }\ell,j.
\]
这些等式适用于非交换环，系数均在左模元素的左侧，乘积的次序是 \(a_{\ell i}b_{ij}\)。
\end{corollary}
\begin{proof}
平坦性给出上一定理的单关系分解。上一定理证明中对关系行数的归纳，把它推广为所述有限关系组分解。
\end{proof}

这里的同时分解比逐条分别分解强。若用 \(A_\ell=(a_{\ell i})_i\) 表示第 \(\ell\) 个系数行，逐条分解只给出 \(\forall\ell\,\exists B_\ell,Y_\ell\)，使 \(X=B_\ell Y_\ell\)、\(A_\ell B_\ell=0\)；同时分解则给出 \(\exists B,Y\,\forall\ell\)，使 \(X=BY\)、\(A_\ell B=0\)。后一种要求让全部关系共用同一个系数矩阵和元素族，才能在有限展示中同时保持全部展示关系。

\ProseParagraphBreak
在交换环上，一条关系的有限多个系数生成一个有限生成理想 \(I=(a_1,\ldots,a_n)\)。它们所对应的张量 \(\sum_i a_i\otimes x_i\in I\otimes_RM\) 在乘法下为零。因此，有限关系判据使我们可以把所有右模单射的检验缩到这些理想的包含 \(I\hookrightarrow R\)。

\begin{theorem}[有限生成理想张量判据]\label{b2:thm:flat-ideal-criterion}
设 \(R\) 为交换含幺环，\(M\) 为幺性 \(R\) 模，则 \(M\) 平坦，当且仅当对每个有限生成理想 \(I\subseteq R\)，由包含 \(I\hookrightarrow R\) 诱导的映射
\[
\mu_I:I\otimes_R M\longrightarrow M,\qquad a\otimes m\longmapsto am
\]
均为单射。
\end{theorem}
\begin{proof}
必要性来自 \(I\hookrightarrow R\) 与 \(R\otimes_R M\cong M\)。为证充分性，给定关系 \(\sum_i a_i x_i=0\)，令 \(I=(a_1,\ldots,a_n)\)。元素 \(\sum_i a_i\otimes x_i\) 经 \(\mu_I\) 映到零，故由假设在 \(I\otimes_R M\) 中已为零。

取满射 \(q:R^n\to I\)、\(e_i\mapsto a_i\)，令 \(K=\ker q\)。张量右正合性说明 \((x_i)\in R^n\otimes_R M\) 来自 \(K\otimes_R M\)。写出一个原像 \(\sum_j k_j\otimes y_j\)，再展开 \(k_j\) 的坐标，便得到 \eqref{b2:eq:flat-relation-factorization}。因此有限关系判据适用，\(M\) 平坦。
\end{proof}

这个理想判据不要求 \(R\) 是 Noether 环\footnote{诺特（Emmy Noether，1882-1935），德国数学家，推进了抽象代数中的理想理论，并发现对称性与守恒律之间的联系。}。有限性发生在一条关系所涉及的系数表中，所以只检验有限生成理想已经足够；不要把它误读成“所有理想本来都有限生成”。本节末的\cref{b2:prop:dual-number-flat}将用双数环的理想检验，把平坦性与一个幂零算子的核、像及 Jordan 块联系起来。

\subsection{局部化的平坦性}
\label{b2:subsec:0703:05}

设 \(R\) 是交换含幺环，\(S\subseteq R\) 是含 \(1\) 的乘法封闭集。由\cref{b2:prop:localization-exact,b2:prop:localization-tensor}，模的局部化保持短正合列，并由 \(S^{-1}R\otimes_R-\) 实现；交换性又允许交换两个张量因子，因而 \(S^{-1}R\) 的平坦性已经成立。有限关系判据还给出更具体的解释：分式中的有限关系，可以化为原环中的真正零关系，再把分母集中到一个共同元素中。

\begin{proposition}\label{b2:prop:localization-flat}
设 \(R\) 是交换含幺环，\(S\subseteq R\) 是含 \(1\) 的乘法封闭集，则局部化 \(S^{-1}R\) 作为幺性 \(R\) 模平坦。
\end{proposition}
\begin{proof}
设 \(\sum_i a_i x_i=0\) 是局部化中的有限关系。将有限多个 \(x_i\) 的分母取乘积，通分为 \(x_i=r_i/s\)，其中 \(s\in S\)。由\eqref{b2:eq:localization-zero}，存在 \(u\in S\) 使 \(u\sum_i a_i r_i=0\)。这使 \(b_i=ur_i\) 在原环中满足 \(\sum_i a_i b_i=0\)。为了同时恢复每个 \(x_i\)，把共同分母和这个额外因子一起放入 \(y=1/(us)\)，便有
\[
x_i=\frac{r_i}s=(ur_i)\frac1{us}=b_i y.
\]
于是所有分式由同一个 \(y\) 重构，而系数 \(b_i\) 已在 \(R\) 中满足真正的零关系。这正是\cref{b2:thm:flat-finite-relations}所需的分解。若 \(0\in S\)，局部化为零模，仍满足同一论证或直接由定义得到平坦性。
\end{proof}

第三卷第12.1节将进一步证明如何在素理想或极大理想处检测一般模的平坦性；本节的\cref{b2:thm:flat-ideal-criterion}是其中使用的一项已证工具。第三卷第13.4节再讨论 Noether 条件下的完备化。完备化涉及逆极限，不能仅因局部化平坦，就断言完备化也平坦；相关有限性与正合性论证留在其主要证明位置。

\ProseParagraphBreak
在有非零幂零元的环上，乘一个标量本身未必是单射，理想张量判据则仍能直接应用。双数环把这一检验化为一个平方为零的算子的核像条件。

\begin{proposition}\label{b2:prop:dual-number-flat}
设 \(k\) 是域，\(R=k[t]/(t^2)\)，记 \(\varepsilon=t+(t^2)\)。对任意幺性 \(R\) 模 \(M\)，令 \(E:M\to M\)、\(E(m)=\varepsilon m\)，则
\[
M\text{ 平坦}\quad\Longleftrightarrow\quad\ker E=\varepsilon M.
\]
若 \(M\) 作为 \(k\) 向量空间有限维，则这两个条件还等价于 \(E\) 的 Jordan 块全为 \(J_2(0)\)，也等价于存在非负整数 \(r\) 使 \(M\cong R^r\)。零模对应 \(r=0\) 及空的 Jordan 块族。
\end{proposition}
\begin{proof}
环 \(R\) 的理想只有 \(0,(\varepsilon),R\)。事实上，每个元素唯一写成 \(a+b\varepsilon\)，\(a,b\in k\)；若 \(a\ne0\)，它的逆为 \(a^{-1}-a^{-2}b\varepsilon\)，所以包含它的理想等于 \(R\)。一个真理想的元素因而都在一维 \(k\) 空间 \(k\varepsilon\) 中，作为 \(k\) 子空间只能为零或整个 \(k\varepsilon\)。这些理想均有限生成，且理想 \(0\) 与 \(R\) 的张量映射显然单射。因此\cref{b2:thm:flat-ideal-criterion}把平坦性化为检验 \((\varepsilon)\otimes_RM\to M\) 的单射性。

映射 \(R\to(\varepsilon)\)、\(a\mapsto a\varepsilon\) 满射，核为 \((\varepsilon)\)，故给出 \(R\) 模同构 \(R/(\varepsilon)\cong(\varepsilon)\)。由\cref{b2:prop:tensor-quotient}，
\[
(\varepsilon)\otimes_RM\cong M/\varepsilon M,
\qquad \varepsilon\otimes m\longleftrightarrow m+\varepsilon M.
\]
理想张量映射在这个同构下为 \(m+\varepsilon M\mapsto\varepsilon m\)。因 \(E^2=0\)，总有 \(\varepsilon M\subseteq\ker E\)，而该映射的核正是 \((\ker E)/\varepsilon M\)。它单射当且仅当 \(\ker E=\varepsilon M\)，这证明任意 \(M\) 的判据。

若 \(M\) 有限 \(k\) 维，\(E^2=0\) 使其最小多项式整除 \(t^2\)，由\cref{b2:thm:jordan-canonical-form}，非零 \(M\) 上的块只能为 \(J_1(0)\) 或 \(J_2(0)\)。在一个 \(J_1(0)\) 块上，核是一维空间而像为零；在一个 \(J_2(0)\) 块上，核与像都是同一个一维空间。核、像按块取直和，故 \(\ker E=\varepsilon M\) 恰好排除所有 \(J_1(0)\) 块。

每个 \(J_2(0)\) 块可取基 \(\varepsilon v,v\)，其中 \(\varepsilon v\ne0\)。映射 \(R\to k\varepsilon v+kv\)、\(a+b\varepsilon\mapsto av+b\varepsilon v\) 是 \(R\) 线性双射，所以 \(r\) 个这样的块给出 \(M\cong R^r\)。反过来，正则模 \(R\) 中 \(\ker E=k\varepsilon=\varepsilon R\)，有限份直和同样满足核像条件；也可由\cref{b2:thm:projective-flat}直接得到自由模的平坦性。零模满足全部条件。
\end{proof}

% !TeX root = ../../Book2.tex
% 纲要标题：### 7.4 限制条件与反例
% 纲要位置：计划纲要.md，第 1295 行。

\section{限制条件与反例}
\label{b2:sec:0704}

自由、投射与平坦之间的蕴含已经建立，但逆命题要逐一检查环及模的条件。\(\mathbb Q\) 已经给出平坦而不投射的例子；本节先说明无挠为什么更弱，再说明有限展示如何补足从平坦到投射所缺少的条件。

% 纲要要点（供目录核对）：
%
% * 一般整环上无挠不等于平坦
% * 有限展示平坦模的投射性与左 Noether 环上的有限生成推论
% * Ext、Tor 对投射、内射、平坦模的刻画；构造与证明在第四卷第8、9章
%
% ---
%

\subsection{无挠与平坦的区别}
\label{b2:subsec:0704:01}

取域 \(k\) 上的多项式环 \(R=k[x,y]\)、\(I=(x,y)\)。由于 \(I\subseteq R\)，它无挠，且由两个元素生成。无挠性排除的是一元关系 \(am=0\)（\(a\ne0\)）中的非零 \(m\)；这里两个生成元却有关系 \((-y)x+xy=0\)。下面将看到，这条二元关系留下的非零张量会被包含 \(I\hookrightarrow R\) 消去。

\begin{example}\label{b2:exm:torsionfree-not-flat}
设 \(k\) 为域，\(R=k[x,y]\)。理想 \(I=(x,y)\) 是有限展示无挠 \(R\) 模，却不是平坦模。
\end{example}
\begin{proof}
满射 \(q:R^2\to I\)、\((a,b)\mapsto ax+by\) 的核由 \((-y,x)\) 生成。事实上，若 \(ax+by=0\)，在 \(R/(x)=k[y]\) 中得到 \(\bar b y=0\)，所以 \(b=xc\)；再由整环中可消去 \(x\)，得到 \(a=-yc\)。反向包含直接成立。因此
\[
R\xrightarrow{c\mapsto(-yc,xc)}R^2\xrightarrow{q}I\longrightarrow0
\]
是一个有限展示。

与 \(I\) 张量，利用\cref{b2:cor:tensor-finite-presentation}对关系矩阵 \(\binom{-y}{x}\) 的计算，得到
\[
I\otimes_R I\cong (I\oplus I)/\bigl\{\,(-yc,xc):c\in I\,\bigr\}.
\]
在此识别下，\(x\otimes y-y\otimes x\) 对应 \((y,-x)\) 的陪集。若它为零，就有 \(c\in I\) 使 \(y=-yc\)、\(-x=xc\)；消去非零的 \(x\) 得 \(c=-1\)，与 \(I\) 为真理想矛盾。因此这个张量非零。

然而包含 \(I\hookrightarrow R\) 张量 \(I\) 后，把它送到 \(R\otimes_R I\cong I\) 中的 \(xy-yx=0\)。这给出未被保持的单射，故 \(I\) 不平坦。
\end{proof}

张量 \(x\otimes y-y\otimes x\) 在乘法下成为零，却未由张量积本身的关系化为零。这个障碍涉及两个系数共同作用。无挠性只控制一个非零标量的零化关系，平坦性则要控制任意有限关系。

\subsection{有限展示平坦模的投射性}
\label{b2:subsec:0704:02}

若要将模的生成元送入自由模，必须保证这些像满足原模的全部关系。平坦性的有限关系判据能够同时处理一个有限关系系统；有限展示恰好保证全部关系由这样的系统生成，因而能够构造一个定义在整个商模上的同态。

\begin{theorem}\label{b2:thm:fp-flat-projective}
设 \(R\) 是含幺环，\(M\) 是左 \(R\) 模。\(M\) 有限展示且平坦，当且仅当它有限生成且投射。
\end{theorem}
\begin{proof}
设 \(M\) 有限展示，取生成元 \(x_1,\ldots,x_n\)，以及生成全部关系的有限个关系
\[
\sum_{i=1}^n a_{\ell i}x_i=0\qquad(1\le\ell\le s).
\]
\(M\) 又平坦，故可用\cref{b2:cor:flat-finite-relation-systems}，找到 \(y_1,\ldots,y_t\) 与 \(b_{ij}\)，使
\[
x_i=\sum_jb_{ij}y_j,\qquad
\sum_i a_{\ell i}b_{ij}=0\quad\text{对所有}\;\ell,j.
\]
记 \(A=(a_{\ell i})\)、\(B=(b_{ij})\)、\(X=(x_i)^{\mathsf T}\)、\(Y=(y_j)^{\mathsf T}\)。上述等式就是 \(X=BY\)、\(AB=0\)，其中乘积始终按原次序计算。在左自由模 \(R^t\) 中记标准基为 \(e_j\)。以 \(B\) 的第 \(i\) 行规定 \(x_i\mapsto\sum_j b_{ij}e_j\)，则第 \(\ell\) 条关系的像为 \(\sum_j(\sum_i a_{\ell i}b_{ij})e_j=0\)。因此这个赋值杀掉所有展示关系及其左线性组合，诱导左模同态 \(h:M\to R^t\)。再用 \(Y\) 定义 \(g:R^t\to M\)、\(e_j\mapsto y_j\)。等式 \(X=BY\) 给出 \(gh(x_i)=x_i\)，所以 \(gh=\operatorname{id}_M\)。由\cref{b2:prop:finite-projective-summand}，\(M\) 有限生成且投射。

反过来，有限生成投射模 \(M\) 是某个 \(R^n\) 的直和因子，写 \(R^n=M\oplus K\)。\(K\) 也是 \(R^n\) 的同态像，因此有限生成。选有限自由满射 \(R^s\to K\)，再复合包含 \(K\to R^n\)，得到有限展示 \(R^s\to R^n\to M\to0\)。投射模平坦则是\cref{b2:thm:projective-flat}。
\end{proof}

若环的左理想都有限生成，即环是左 Noether 环，有限生成模的关系也能由有限多个关系生成。于是得到下面的推论。

\begin{corollary}\label{b2:cor:noetherian-fg-flat-projective}
设 \(R\) 为左 Noether 环。每个有限生成平坦左 \(R\) 模都是投射模。
\end{corollary}
\begin{proof}
先对 \(n\) 归纳证明 \(R^n\) 的每个左子模有限生成。\(n=1\) 时这就是左理想的假设。对于 \(K\subseteq R^n\)，最后坐标投影的像是有限生成左理想，核可视为 \(R^{n-1}\) 的子模，由归纳假设有限生成。把像的有限生成元提升到 \(K\)，再添上核的有限生成元，就生成 \(K\)。

对有限生成左模 \(M\) 取满射 \(R^n\to M\)，其核因上述结论而有限生成。因此 \(M\) 有限展示；若它又平坦，\cref{b2:thm:fp-flat-projective}便给出投射性。
\end{proof}

以下例子说明，一般环上有限生成并不足以保证有限展示。

\begin{example}[有限生成平坦但非投射的循环模]\label{b2:exm:fg-flat-nonprojective}
设 \(k\) 是域，\(R=\prod_{n\ge1}k\)，\(I=\bigoplus_{n\ge1}k\) 是有限支集序列组成的理想。循环模 \(M=R/I\) 平坦，却不投射；它也不有限展示。
\end{example}
\begin{proof}
给定有限关系 \(\sum_{i=1}^r a_i\bar b_i=0\)，选取 \(b_i\in R\) 作代表，令 \(c=\sum_i a_i b_i\in I\)。这条关系在商模中成立，在环中却可能有误差 \(c\)。误差有有限支集，故可取有限坐标集 \(F\) 包含它的支集，并令 \(e_F\) 是这些坐标为一、其余为零的幂等元。它满足 \(e_Fc=c\)，相当于只对这项误差起单位作用；乘 \(1-e_F\) 就消去了误差。因为 \(e_Fb_i\in I\)，这一修正不改变 \(\bar b_i\)，于是有 \(\bar b_i=(1-e_F)b_i\bar1\)，而
\[
\sum_i a_i(1-e_F)b_i=(1-e_F)c=0.
\]
所有 \(\bar b_i\) 由共同的元素 \(\bar1\) 重构，且其系数 \((1-e_F)b_i\) 在环中满足真正的零关系。由\cref{b2:thm:flat-finite-relations}，\(M\) 平坦。

若 \(M\) 投射，自然满射 \(R\to M\) 分裂，因而 \(I\) 是正则模的直和因子。相应投影 \(p:R\to I\) 由 \(e=p(1)\in I\) 决定，其像为 \(Re\)，所以 \(I=Re\)。但 \(e\) 只有有限支集，\(Re\) 无法包含支集在其外的单坐标向量，矛盾。相同的支集论证说明 \(I\) 不是有限生成理想，由\cref{b2:prop:finite-presentation-kernel}，\(R/I\) 不有限展示。它只有一个生成元 \(\bar1\)，但这个生成元的全部关系组成的理想 \(I\) 不能由有限多条关系生成。
\end{proof}

\subsection{Ext、Tor 刻画}
\label{b2:subsec:0704:03}

提升、延拓和张量正合性还可以用记录其失败的交换群来研究。设 \(R\) 为含幺环，\(P,E,M\) 为幺性左 \(R\) 模。第四卷将构造导出函子 \(\operatorname{Ext}\) 与 \(\operatorname{Tor}\)，并证明下列刻画；这里先说明公式与本章定义的联系：
\[
\begin{aligned}
P\;\text{投射}&\quad\Longleftrightarrow\quad
\operatorname{Ext}_R^1(P,N)=0\ \text{对每个左模}\;N,\\
E\;\text{内射}&\quad\Longleftrightarrow\quad
\operatorname{Ext}_R^1(N,E)=0\ \text{对每个左模}\;N,\\
M\;\text{平坦}&\quad\Longleftrightarrow\quad
\operatorname{Tor}_1^R(N,M)=0\ \text{对每个右模}\;N.
\end{aligned}
\]
消解及 Ext、Tor 的构造见第四卷第8、9章，具体定义在第9.2—9.3节，上述刻画的证明在第9.4节。尤其第三行须让右模 \(N\) 位于张量积的左侧，与本章平坦左模的约定一致。
\symidx{\(\operatorname{Ext}_R^1(P,N)\)：记录模扩张障碍的一次 Ext 群}
\symidx{\(\operatorname{Tor}_1^R(N,M)\)：记录张量正合性障碍的一次 Tor 群}

\ProseParagraphBreak
第一行说，以 \(P\) 为商模的短正合列都能分裂，相应的 \(\operatorname{Ext}^1\) 不留下扩张障碍；第二行说，以 \(E\) 为子模的短正合列都能分裂。第三行则说，任意右模与 \(M\) 配对时，\(\operatorname{Tor}_1\) 不留下张量正合性的障碍。本章已经用提升、延拓和单射保持完整刻画了这些性质；上述符号的定义、计算、消解选择无关性和长正合列将在第四卷建立。

\ProseParagraphBreak
\cref{b2:tab:projective-injective-flat-exactness}将本章的三个定义并列。前两行的 \(A,B,C\) 属于左 \(R\) 模短正合列，最后一行属于右 \(R\) 模短正合列；\(P,E,M\) 均为左 \(R\) 模。各性质所补足的正是相应函子原先未必保持的那个箭头。
\begin{table}[htbp]
\centering
\caption{投射、内射与平坦的正合性比较}\label{b2:tab:projective-injective-flat-exactness}
\begin{tabularx}{\linewidth}{@{}p{0.13\linewidth}p{0.31\linewidth}X@{}}
\toprule
性质 & 函子 & 补足的正合性及映射任务 \\
\midrule
投射 \(P\) & \(\operatorname{Hom}_R(P,-)\) & \(\operatorname{Hom}_R(P,B)\to\operatorname{Hom}_R(P,C)\) 满射，即沿 \(B\twoheadrightarrow C\) 提升同态。 \\
内射 \(E\) & \(\operatorname{Hom}_R(-,E)\) & \(\operatorname{Hom}_R(B,E)\to\operatorname{Hom}_R(A,E)\) 满射，即沿 \(A\hookrightarrow B\) 延拓同态。 \\
平坦 \(M\) & \(-\otimes_R M\) & \(A\otimes_R M\to B\otimes_R M\) 单射，即张量保持单射。 \\
\bottomrule
\end{tabularx}
\end{table}


\begin{chaptersummary}
投射模是自由模的直和因子，其 Hom 函子保持满射；内射模允许延拓到任意较大的定义域，其反变 Hom 函子保持相应的正合列。Baer 判据把延拓问题缩减到左理想。可除交换群提供基本的内射对象，再由 \(\mathbb Q/\mathbb Z\) 分离元素和余诱导构造，得到任意环上足够多的内射模。

\ProseParagraphBreak
平坦性补足张量积缺少的左端单射。有限关系判据把这个函子条件变为系数的因子分解，交换环上再化为有限生成理想的张量检验。主理想整环上，无挠恰好等价于平坦；一般整环上，理想 \((x,y)\) 表明两者可以不同。有限展示平坦模的全部关系能够同时处理，从而可写成有限自由模的直和因子。\(\mathbb Q\) 及无限积环的循环商模分别说明，平坦本身和有限生成平坦都不足以推出投射。

\ProseParagraphBreak
对任意含幺环，总有
\[
\text{自由}\ \Longrightarrow\ \text{投射}\ \Longrightarrow\ \text{平坦}.
\]
整环上平坦还蕴含无挠；主理想整环上平坦与无挠等价。若模另有有限生成条件，结构定理又将无挠模化为有限自由模，因而这四个条件全部等价。在左 Noether 环上，有限生成平坦左模也投射，因为有限生成保证有限展示。各处的环和有限性假设不能省去；内射性与上述链在一般环上没有横向蕴含。
\end{chaptersummary}

\BookExercises

\begin{exercise}\label{b2:ex:07-projective-quotient-idempotent}
设 \(R\) 为交换环，\(I\) 为理想。证明 \(R/I\) 是投射 \(R\) 模，当且仅当 \(I=Re\)，其中 \(e^2=e\)。据此判断 \(\mathbb Z/6\mathbb Z\) 作为 \(\mathbb Z\) 模和作为 \(\mathbb Z/12\mathbb Z\) 模是否投射。
\end{exercise}
\begin{solution}
若 \(R/I\) 投射，\(R\to R/I\) 分裂，故存在投影 \(p:R\to I\)，在 \(I\) 上为恒等。令 \(e=p(1)\)，则 \(p(r)=re\)，因而 \(I=Re\)；又 \(p(e)=e\) 给出 \(e^2=e\)。反过来，\(R=Re\oplus R(1-e)\)，所以 \(R/Re\cong R(1-e)\) 投射。

整数环的幂等元只有零和一，\(6\mathbb Z\) 不是相应理想，所以第一种情况不投射。第二种情况的核为 \((\bar6)\subseteq\mathbb Z/12\mathbb Z\)。此理想只有 \(\bar0,\bar6\)，而 \(\bar6^2=\bar0\ne\bar6\)，不存在生成它的幂等元，所以也不投射。作为对照，\(\mathbb Z/3\mathbb Z\) 是 \(\mathbb Z/6\mathbb Z\) 上的投射模：核 \((\bar3)\) 由幂等元 \(\bar3\) 生成。
\end{solution}

\begin{exercise}\label{b2:ex:07-explicit-lifting}
设 \(P\) 有有限对偶基 \((p_i,\varphi_i)\)，\(q:B\to C\) 满射，\(f:P\to C\)。选取 \(b_i\in B\) 使 \(q(b_i)=f(p_i)\)，写出一个提升的公式。若 \(g_0\) 是一个提升，描述所有提升与 \(\operatorname{Hom}_R(P,\ker q)\) 的关系。
\end{exercise}
\begin{solution}
令 \(g(p)=\sum_i\varphi_i(p)b_i\)。每个 \(\varphi_i\) 左线性，故 \(g\) 左线性，而且
\[
qg(p)=\sum_i\varphi_i(p)f(p_i)
=f\left(\sum_i\varphi_i(p)p_i\right)=f(p).
\]
任意另一提升 \(g\) 满足 \(q(g-g_0)=0\)，故 \(g-g_0\) 唯一经 \(\ker q\) 的包含分解。反之，给 \(g_0\) 加上任意从 \(P\) 到 \(\ker q\) 的同态，仍为提升。因此选定 \(g_0\) 后，所有提升与该 Hom 群一一对应；没有指定 \(g_0\) 时，不应把某一个提升称为天然的零元。
\end{solution}

\begin{exercise}\label{b2:ex:07-idempotent-matrix}
在整数模 \(\mathbb Z^2\) 上取 \(E=\begin{pmatrix}2&-2\\1&-1\end{pmatrix}\)。计算 \(E^2\)、\(\operatorname{im}E\)、\(\ker E\)，并求一个 \(U\in\operatorname{GL}_2(\mathbb Z)\)，使 \(U^{-1}EU=\operatorname{diag}(1,0)\)。对 \(P=\operatorname{im}E\)，分别写出一组只有一项的有限对偶基，以及由 \(Ee_1,Ee_2\) 和原坐标泛函得到的两项有限对偶基。比较两者的系数是否唯一。
\end{exercise}
\begin{solution}
令 \(u=(2,1)^{\mathsf T}\)、\(w=(1,1)^{\mathsf T}\)，则 \(E(a,b)^{\mathsf T}=(a-b)u\)，且 \(Eu=u\)、\(Ew=0\)。因此 \(E^2=E\)，\(P=\mathbb Zu\)、\(\ker E=\mathbb Zw\)。取 \(U=(u\mid w)=\begin{pmatrix}2&1\\1&1\end{pmatrix}\)，其行列式为一，且 \(EU=U\operatorname{diag}(1,0)\)，即得所需换基。

泛函 \(\lambda:P\to\mathbb Z\)、\(\lambda(au)=a\) 与 \(u\) 构成一项有限对偶基，此时 \(\lambda(u)=1\)，也是通常的基及坐标泛函。另一组取 \(p_1=u\)、\(p_2=-u\)，原坐标限制为 \(\varphi_1(au)=2a\)、\(\varphi_2(au)=a\)。于是 \(\varphi_1(au)p_1+\varphi_2(au)p_2=au\)。后一组给出指定的重构系数，但 \(cp_1+dp_2=(c-d)u\)，所以表达系数并不唯一，且 \(\varphi_i(p_j)\) 也不是 \(\delta_{ij}\)。
\end{solution}

\begin{exercise}\label{b2:ex:07-pid-finite-projective}
设 \(R\) 为主理想整环，\(E\in M_n(R)\) 满足 \(E^2=E\)。证明存在 \(U\in\operatorname{GL}_n(R)\) 及 \(0\le r\le n\)，使 \(U^{-1}EU=\operatorname{diag}(I_r,0)\)。若 \(R=k\times k\)，取一阶矩阵 \(E=(e)\)、\(e=(1,0)\)，说明这一结论为何失败，并指出证明中使用了主理想整环的哪项性质。
\end{exercise}
\begin{solution}
幂等性给出 \(R^n=\operatorname{im}E\oplus\ker E\)。这两个直和因子都是有限生成无挠模，故由\cref{b2:thm:pid-module-structure}自由。分别取基 \(u_1,\ldots,u_r\) 和 \(v_1,\ldots,v_s\)；两组基合起来是 \(R^n\) 的基，因秩不变性有 \(r+s=n\)。以它们为列构成的 \(U\) 可逆。\(E\) 在第一组基上为恒等，在第二组基上为零，所以得到所述相似形式。

在 \(R=k\times k\) 中，一阶矩阵的相似变换不改变 \(e\)，因为 \(R\) 交换；而 \(e\) 既非零也非一，不可能变为 \(\operatorname{diag}(I_r,0)\)（此时 \(r=0\) 或一）。主理想整环条件保证有限生成无挠模自由，故两直和因子能选出通常的基；一般环上的投射直和因子不一定自由。
\end{solution}

\begin{exercise}\label{b2:ex:07-baer-pid}
设 \(R\) 为主理想整环，\(D\) 是 \(R\) 模。若对每个 \(0\ne a\in R\)，乘法映射 \(D\xrightarrow{\times a}D\) 都满射，就称 \(D\) 可除。用 Baer 判据证明 \(D\) 内射当且仅当可除。令 \(K\) 为 \(R\) 的分式域，判断 \(K\) 与 \(K/R\) 的内射性。
\end{exercise}
\begin{solution}
若 \(D\) 内射，对 \(d\in D\) 定义 \(aR\to D\)、\(ar\mapsto rd\)。因 \(R\) 为整环，此定义与表示无关。延拓到 \(R\)，令 \(x\) 为一的像，则 \(ax=d\)。反过来，每个非零理想为 \(aR\)。给定 \(f:aR\to D\)，取 \(x\) 使 \(ax=f(a)\)，则 \(r\mapsto rx\) 延拓 \(f\)；零理想取零延拓。因此 Baer 判据成立。\(K\) 与 \(K/R\) 中均可将分式代表元除以 \(a\)，故两者可除并内射。
\end{solution}

\begin{exercise}\label{b2:ex:07-divisible-exact-sequence}
设 \(D\) 为可除交换群，\(n>1\)。给定交换群的短正合列
\(0\to D\xrightarrow{i}B\xrightarrow{q}\mathbb Z/n\mathbb Z\to0\)，选取 \(b\in B\) 使 \(q(b)=\bar1\)，直接修正 \(b\) 以构造 \(q\) 的截面，并描述全部截面的差别。再证明
\[
0\longrightarrow C_n\longrightarrow\mathbb Q/\mathbb Z
\xrightarrow{\times n}\mathbb Q/\mathbb Z\longrightarrow0,
\qquad C_n=\langle1/n+\mathbb Z\rangle,
\]
不分裂。两项结论如何体现内射对象在短正合列中的位置？
\end{exercise}
\begin{solution}
因 \(nb\in\ker q=i(D)\)，写 \(nb=i(d)\)。在 \(D\) 中取 \(a\) 使 \(na=d\)，则 \(b'=b-i(a)\) 满足 \(nb'=0\)、\(q(b')=\bar1\)，故 \(\bar r\mapsto rb'\) 定义截面。选定一个截面后，其他截面与它的差是到 \(D\) 的同态，由 \(\bar1\) 的像决定，且此像属于 \(D[n]=\{d\in D:nd=0\}\)。反过来，每个这样的像都给出一个新截面。

第二列的核是 \(C_n\)，乘 \(n\) 因可除性满射，所以确实正合。若它分裂，就有缩回 \(\mathbb Q/\mathbb Z\to C_n\)。但任意这样的同态 \(f\) 都为零：对任意 \(z\) 取 \(z'\) 使 \(nz'=z\)，则 \(f(z)=nf(z')=0\)，因为 \(nC_n=0\)。这与缩回在非零的 \(C_n\) 上为恒等矛盾。

第一列中左端 \(D\) 内射，故能分裂，显式修正也给出了截面；第二列的中间项和右端虽都内射，左端却不是，因而不能据内射性断言分裂。
\end{solution}

\begin{exercise}\label{b2:ex:07-finite-characters}
设 \(A\) 是有限交换群，记 \(A^\vee=\operatorname{Hom}_{\mathbb Z}(A,\mathbb Q/\mathbb Z)\)。证明 \(|A^\vee|=|A|\)，并证明自然求值映射 \(A\to(A^\vee)^\vee\) 是同构。
\end{exercise}
\begin{solution}
对 \(A=\mathbb Z/n\mathbb Z\)，同态由 \(\bar1\) 的像确定，其像可取 \(j/n+\mathbb Z\)、\(0\le j<n\)，所以恰有 \(n\) 个。有限交换群由\cref{b2:thm:pid-module-structure}分成有限个循环群的直和；Hom 把有限直和变成有限直积，因而个数相乘，得到 \(|A^\vee|=|A|\)。

求值映射把 \(a\) 送到 \(\chi\mapsto\chi(a)\)，为群同态。\cref{b2:lem:qz-separates}说明非零 \(a\) 总被某个 \(\chi\) 检测到，故它单射。两端都是同阶有限群，因而同构。映射的公式不需选循环分解，分解只用于证明两端阶数相同。
\end{solution}

\begin{exercise}\label{b2:ex:07-injective-summands}
证明内射模的直和因子内射，并给出内射却不平坦的整数模。再说明为什么“一个模嵌入内射模”不能推出它内射。
\end{exercise}
\begin{solution}
若 \(E\) 是内射模 \(J\) 的直和因子，记包含与投影为 \(i,p\)，且 \(pi=\operatorname{id}_E\)。对子模上的 \(f:A\to E\)，先延拓 \(if\) 到较大模 \(B\)，再与 \(p\) 复合，即为 \(f\) 的延拓。

\(\mathbb Q/\mathbb Z\) 可除而内射，但 \(1/2+\mathbb Z\) 为非零挠元，故由\cref{b2:prop:flat-torsionfree}不平坦。\(\mathbb Z\) 嵌入内射模 \(\mathbb Q\)，却不是内射；这也说明上述直和因子的缩回条件不可省去。
\end{solution}

\begin{exercise}\label{b2:ex:07-group-algebra-injective}
设 \(G\) 为有限群，\(k\) 为任意域，\(R=k[G]\)。在 \(J=\operatorname{Hom}_k(R,k)\) 上取左作用 \((rf)(s)=f(sr)\)。证明 \(J\) 是内射左 \(R\) 模，并用群元基证明 \(R\cong J\) 为左模，从而正则左模 \(R\) 内射。这里不假设 \(\operatorname{char}k\) 与 \(|G|\) 互素。
\end{exercise}
\begin{solution}
与\cref{b2:lem:coinduced-injective}相同的显式公式给出 \(\operatorname{Hom}_R(M,J)\cong\operatorname{Hom}_k(M,k)\)。向量空间的子空间有补空间，故任意 \(k\) 线性泛函可延拓；利用此对应，\(J\) 内射。

令 \(\lambda:R\to k\) 取单位群元的系数，定义 \(\Phi(a)(s)=\lambda(sa)\)。有 \(\Phi(ra)(s)=\lambda(sra)=\bigl(r\cdot\Phi(a)\bigr)(s)\)，故 \(\Phi\) 左线性。若 \(a=\sum_g a_g g\)，则 \(\Phi(a)(g^{-1})=a_g\)，因而 \(\Phi\) 单射。两端维数均为 \(|G|\)，所以同构。这项证明未对群阶作除法，因此对任意特征成立。

内射性来自 \(R\cong\operatorname{Hom}_k(R,k)\) 的这种自对偶结构；这是群代数的 Frobenius 型性质，即使特征整除 \(|G|\) 仍成立。因此这里不能将正则模的内射性归因于 Maschke 定理中的半单性。
\end{solution}

\begin{exercise}\label{b2:ex:07-integer-relation}
设 \(M\) 为无挠交换群，且 \(6x+10y=0\)。求 \(z\in M\) 使 \(x=5z\)、\(y=-3z\)，从而给出这条关系的一项有限关系分解。去掉无挠条件后结论是否仍成立？
\end{exercise}
\begin{solution}
由无挠性可将关系约去二，得到 \(3x+5y=0\)。取 \(z=2x+3y\)，则 \(5z-x=9x+15y=0\)，而 \(-3z-y=-6x-10y=0\)。系数满足 \(6\cdot5+10\cdot(-3)=0\)，正是有限关系判据的形式。

在 \(M=\mathbb Z/2\mathbb Z\) 中取 \(x=\bar1,y=\bar0\)，原关系成立，但所求等式分别要求 \(z=\bar1\) 与 \(z=\bar0\)，矛盾。因此无挠条件实际用于约去二。
\end{solution}

\begin{exercise}\label{b2:ex:07-localization-nonprojective}
设 \(p\) 为素数。证明 \(\mathbb Z[1/p]\) 是平坦而非投射的整数模，并判断它是否内射。
\end{exercise}
\begin{solution}
它是整数环的局部化，由\cref{b2:prop:localization-flat}平坦。若 \(h:\mathbb Z[1/p]\to\mathbb Z\)，每个 \(h(x)\) 都被任意 \(p^n\) 整除，因为 \(x/p^n\) 仍在定义域内，所以 \(h=0\)。若该模投射，它是自由模的直和因子，其嵌入与每个整数坐标的复合均为零，导致嵌入为零，矛盾。

另取素数 \(q\ne p\)，方程 \(qx=1\) 在 \(\mathbb Z[1/p]\) 中无解，否则有 \(q a=p^n\) 的整数等式。因此该群不可除，由\cref{b2:thm:divisible-injective}不内射。
\end{solution}

\begin{exercise}\label{b2:ex:07-dual-numbers-flat}
令 \(R=k[\varepsilon]/(\varepsilon^3)\)，\(M\) 为有限维 \(k\) 向量空间上的 \(R\) 模，记乘 \(\varepsilon\) 的算子为 \(T\)。证明下列条件等价：\(M\) 平坦；\(\ker T^2=\operatorname{im}T\)；\(T\) 的 Jordan 块全部为 \(J_3(0)\)；\(M\) 为自由 \(R\) 模。计算理想 \((\varepsilon)\) 和 \((\varepsilon^2)\) 的张量检验各对应哪个核像条件，并说明为什么这里检验第二个理想就已足够。
\end{exercise}
\begin{solution}
因 \(T^3=0\)，Jordan 块只能为 \(J_1(0),J_2(0),J_3(0)\)。环的理想是 \(0,(\varepsilon^2),(\varepsilon),R\)：非零元素的最低非零次数分别为零、一、二，提出相应的 \(\varepsilon\) 幂后，剩余因子为单位，便生成相应理想。

同构 \((\varepsilon)\cong R/(\varepsilon^2)\) 给出张量映射 \(M/\varepsilon^2M\to M\)、\(m\mapsto\varepsilon m\)，其单射性等价于 \(\ker T=\operatorname{im}T^2\)。同构 \((\varepsilon^2)\cong R/(\varepsilon)\) 则给出 \(M/\varepsilon M\to M\)、\(m\mapsto\varepsilon^2m\)，其单射性等价于 \(\ker T^2=\operatorname{im}T\)。因此平坦性蕴含后一条件。

在大小 \(d=1,2,3\) 的 Jordan 块上，\(\dim\ker T^2=\min(2,d)\)，\(\dim\operatorname{im}T=d-1\)，又总有 \(\operatorname{im}T\subseteq\ker T^2\)。两维数之差分别为一、一、零。因此核像相等恰好排除大小一、二的块。每个大小三的块同构于正则模 \(R\)，所以仅有 \(J_3(0)\) 时 \(M\) 自由并平坦，且第一个理想的核像条件也自动成立。这就证明所有等价关系及第二个理想检验的充分性。
\end{solution}

\begin{exercise}\label{b2:ex:07-ideal-tensor-kernel}
沿用\cref{b2:exm:torsionfree-not-flat}中的 \(R=k[x,y]\)、\(I=(x,y)\)。证明乘法映射 \(I\otimes_R I\to I^2\) 的整个核同构于 \(R/I\)，由 \(x\otimes y-y\otimes x\) 生成。
\end{exercise}
\begin{solution}
该例给出 \(I\otimes_R I\cong(I\oplus I)/L\)，其中 \(L\) 是全部 \((-yc,xc)\)、\(c\in I\) 组成的子模，乘法为 \((a,b)\mapsto xa+yb\)。在 \(R^2\) 中它的全部零关系为 \((-yd,xd)\)、\(d\in R\)，而这两个坐标总属于 \(I\)。因此所求核是以 \(d\in R\) 为参数的这些向量，模去参数 \(d\in I\) 的向量，即 \(R/I\)。符号相反的生成向量 \((y,-x)\) 对应题中的张量，同样生成整个核。特别地，该核被 \(x,y\) 零化；原模 \(I\) 无挠并不保证它的自张量无挠。
\end{solution}

\begin{exercise}\label{b2:ex:07-flat-base-change}
设 \(R\to S\) 为含幺环同态。证明：若左 \(R\) 模 \(M\) 平坦，则 \(S\otimes_R M\) 是平坦左 \(S\) 模；若 \(P\) 为有限生成投射左 \(R\) 模，则 \(S\otimes_R P\) 为有限生成投射左 \(S\) 模。
\end{exercise}
\begin{solution}
右 \(S\) 模包含 \(U\subseteq V\) 限制标量后仍为右 \(R\) 模包含。张量结合与单位同构给出
\[
U\otimes_S(S\otimes_R M)\cong U\otimes_R M,
\]
对 \(V\) 同样成立，并与包含映射相容。由 \(M\) 平坦，右侧诱导映射单射，故左侧也单射。

若 \(P\oplus Q\cong R^n\)，则扩张标量后有 \((S\otimes_R P)\oplus(S\otimes_R Q)\cong S^n\)，因为张量保持直和与同构。因此所求模为有限自由左 \(S\) 模的直和因子。
\end{solution}

\begin{exercise}\label{b2:ex:07-fp-splitting-formula}
在\cref{b2:thm:fp-flat-projective}的证明中，记 \(\pi:R^n\to M\)、\(e_i\mapsto x_i\)。若选择 \(c_{jk}\in R\) 使 \(y_j=\sum_k c_{jk}x_k\)，写出 \(\pi\) 的一个截面，并解释为什么有限关系数不可被“对每条关系分别找一个分解”代替。
\end{exercise}
\begin{solution}
记证明中构造的 \(h:M\to R^t\) 满足 \(h(x_i)=\sum_jb_{ij}e'_j\)。定义 \(u:R^t\to R^n\)、\(u(e'_j)=\sum_kc_{jk}e_k\)，则 \(\pi u(e'_j)=y_j\)，所以 \(s=uh\) 满足
\[
s(x_i)=\sum_{j,k}b_{ij}c_{jk}e_k,\qquad \pi s(x_i)=x_i.
\]
因此 \(\pi s=\operatorname{id}_M\)。全部关系同时满足 \(AB=0\)，才保证 \(h\) 是定义在商模 \(M\) 上的一个同态。若每条关系分别选择不同的矩阵和中间元素，就没有一个统一赋值可据以定义 \(h\)，也无从拼成同一个截面。
\end{solution}

\begin{exercise}\label{b2:ex:07-flat-quotient-ideal}
设 \(R\) 交换，\(I\) 为理想且 \(R/I\) 平坦。证明 \(I/I^2=0\)。若 \(I\) 有限生成，进一步证明它由一个幂等元生成。解释为什么正文中的无限积环例不与此矛盾。
\end{exercise}
\begin{solution}
将单射 \(I\hookrightarrow R\) 与 \(R/I\) 张量。由\cref{b2:prop:tensor-quotient}，定义域为 \(I/I^2\)，而诱导映射到 \(R/I\) 为零。平坦性使其单射，所以 \(I/I^2=0\)。

商 \(I/I^2\) 保留理想元素模去两项理想元素乘积后的信息；结论说，这些元素全都已经是有限个此类乘积的和，即 \(I=I^2\)。这没有说 \(I=0\)，也没有说其中的生成关系消失。

若 \(I\) 有限生成，则 \(R/I\) 有限展示，由\cref{b2:thm:fp-flat-projective}投射。\cref{b2:ex:07-projective-quotient-idempotent}的证明遂给出 \(I=Re\)、\(e^2=e\)。在无限积环例中，有限支集理想确有 \(I^2=I\)：每个有限支集元素 \(a\) 都满足 \(a=e_Fa\)，其中 \(e_F\in I\) 覆盖其支集。但该理想不有限生成，因而不能应用后一结论。
\end{solution}

\begin{exercise}\label{b2:ex:07-projective-hom-tensor}
设 \(k\) 为域，\(R=k\times k\)、\(e=(1,0)\)，\(N_1,N_2\) 为 \(k\) 向量空间。在 \(N=N_1\times N_2\) 上取 \((a,b)(n_1,n_2)=(an_1,bn_2)\)。对 \(P=Re\)，计算 \(P^*=\operatorname{Hom}_R(P,R)\)、\(N\otimes_RP^*\) 及 \(\operatorname{Hom}_R(P,N)\)，显式写出自然映射
\[
\Theta:N\otimes_R P^*\longrightarrow\operatorname{Hom}_R(P,N),\qquad
n\otimes\varphi\longmapsto\bigl(p\mapsto\varphi(p)n\bigr)
\]
及其逆。再取 \(P'=P\oplus R\)，计算对应的张量模、Hom 模和逆公式。分别沿两坐标投影 \(\pi_i:R\to k\) 扩张标量，求 \(k\otimes_{R,\pi_i}P\) 与 \(k\otimes_{R,\pi_i}P'\) 的维数；这些维数就是两个点处的局部秩。由此判断 \(P,P'\) 是否自由，并解释有限对偶基为何仍能用于这两个模。
\end{exercise}
\begin{solution}
令 \(\lambda:P\hookrightarrow R\) 为包含映射。任意 \(\varphi:P\to R\) 由 \(\varphi(e)\) 决定；因 \(e\varphi(e)=\varphi(e)\)，这个值属于 \(Re\)。反过来，对 \(c\in Re\)，公式 \(\varphi_c(p)=cp\) 定义这样的同态，并满足 \(\varphi_c(e)=c\)。故 \(P^*\cong Re\)，其中 \(c\) 对应 \(c\lambda\)。

同样，\(h:P\to N\) 由 \(h(e)\in eN=N_1\times0\) 决定，逆对应把 \(n\in eN\) 送到 \(h_n(re)=rn\)。因此 \(\operatorname{Hom}_R(P,N)\cong eN\)。张量模也有同构
\[
N\otimes_RP^*\longrightarrow eN,
\qquad n\otimes(c\lambda)\longmapsto cn,
\]
其逆为 \(n\mapsto n\otimes\lambda\)。正向公式平衡，而反向复合使用 \(cn\otimes\lambda=n\otimes(c\lambda)\)，所以这两公式互逆。在这两种识别下，\(\Theta\) 就是 \(eN\) 的恒等映射，因为 \(\Theta(n\otimes c\lambda)(e)=cn\)。故原来的逆公式为 \(h\mapsto h(e)\otimes\lambda\)。这两个模只保留 \(N_1\)，没有来自 \(N_2\) 的分量。

对 \(P'=P\oplus R\)，令 \(\lambda_1(p,r)=p\)、\(\lambda_2(p,r)=r\)。有限直和的 Hom 分解给出
\[
\begin{aligned}
(P')^*&\cong Re\oplus R,\\
N\otimes_R(P')^*&\cong eN\oplus N,\\
\operatorname{Hom}_R(P',N)&\cong eN\oplus N.
\end{aligned}
\]
后一个对应在 \((e,0)\)、\((0,1_R)\) 处取值，所得分量可写为 \(N_1^2\times N_2\)。张量模的对应把 \(n\otimes(c\lambda_1+r\lambda_2)\) 送到 \((cn,rn)\)，所以 \(\Theta\) 在这些坐标下仍是恒等映射。其逆为
\[
h\longmapsto h(e,0)\otimes\lambda_1+h(0,1_R)\otimes\lambda_2.
\]

第一坐标投影把 \(e\) 送到一，第二坐标投影把它送到零。由张量积的单位及商模公式，\(k\otimes_{R,\pi_1}P\cong k\)、\(k\otimes_{R,\pi_2}P=0\)；\(R\) 在两处都变为 \(k\)，故 \(P\) 的局部秩为 \((1,0)\)，\(P'\) 的局部秩为 \((2,1)\)。自由模扩张标量后在两处都有同一个基数的基，因而 \(P,P'\) 均不自由。

尽管如此，\(P\) 有有限对偶基 \((e,\lambda)\)，\(P'\) 有有限对偶基 \(\bigl((e,0),\lambda_1\bigr)\)、\(\bigl((0,1_R),\lambda_2\bigr)\)。例如
\[
(p,r)=\lambda_1(p,r)(e,0)+\lambda_2(p,r)(0,1_R).
\]
第一组的 \(\lambda(e)=e\ne1_R\)，所以它并不是普通基及坐标泛函；有限对偶基要求的是指定的有限线性重构，正是上述逆公式所用的条件。
\end{solution}

\begin{exercise}\label{b2:ex:07-compare-three-classes}
对整数模 \(\mathbb Z\)、\(\mathbb Q\)、\(\mathbb Q/\mathbb Z\)、\(\mathbb Z/n\mathbb Z\)（\(n>1\)），分别判断投射、内射和平坦性。根据答案，指出哪些两类性质之间不存在一般的蕴含。
\end{exercise}
\begin{solution}
\(\mathbb Z\) 自由，所以投射且平坦，但不可除，故不内射。\(\mathbb Q\) 无挠且可除，所以平坦且内射；它到整数的同态全为零，所以不投射。\(\mathbb Q/\mathbb Z\) 可除而有非零挠元，故内射但不平坦，从而也不投射。\(\mathbb Z/n\mathbb Z\) 有挠且不可除，所以不平坦、不内射；投射蕴含平坦，故也不投射。

这些例子说明内射既不蕴含投射，也不蕴含平坦；投射和平坦都不蕴含内射；平坦不蕴含投射。三者之间在任意环上总成立的蕴含是投射推出平坦，其他逆向或横向结论均需另加条件。
\end{solution}
